\documentclass[12pt,a4paper]{article}

% Packages
\usepackage[french]{babel}
\usepackage[utf8]{inputenc}
\usepackage[T1]{fontenc}
\usepackage{graphicx}
\usepackage{geometry}
\usepackage{hyperref}
\usepackage{enumitem}
\usepackage{xcolor}
\usepackage{fancyhdr}
\usepackage{tcolorbox}
\usepackage{listings}
\usepackage{longtable}
\usepackage{array}
\usepackage{booktabs}

% Configuration de la page
\geometry{
    left=2.5cm,
    right=2.5cm,
    top=2.5cm,
    bottom=2.5cm
}

% Configuration des en-têtes et pieds de page
\pagestyle{fancy}
\fancyhf{}
\fancyhead[L]{Guide d'utilisation complet}
\fancyhead[R]{Générateur de Relevés et Attestations}
\fancyfoot[C]{\thepage}

% Configuration des liens
\hypersetup{
    colorlinks=true,
    linkcolor=blue,
    filecolor=magenta,
    urlcolor=cyan,
    pdftitle={Guide d'utilisation - Générateur de Relevés},
    pdfpagemode=FullScreen,
}

% Commandes personnalisées
\newcommand{\menu}[1]{\textbf{\textcolor{blue}{#1}}}
\newcommand{\bouton}[1]{\texttt{\textcolor{teal}{[#1]}}}
\newcommand{\champ}[1]{\textit{\textcolor{purple}{#1}}}
\newcommand{\attention}[1]{\textcolor{red}{\textbf{⚠ #1}}}
\newcommand{\info}[1]{\textcolor{blue}{\textbf{ℹ #1}}}
\newcommand{\astuce}[1]{\textcolor{green}{\textbf{💡 #1}}}

% Début du document
\begin{document}

% Page de titre
\begin{titlepage}
    \centering
    \vspace*{2cm}

    {\Huge\bfseries Guide d'utilisation complet\par}
    \vspace{1cm}
    {\LARGE Générateur de Relevés de Notes\par}
    {\LARGE et Attestations de Réussite\par}
    \vspace{1.5cm}

    % Placeholder pour logo
    \begin{figure}[h]
        \centering
        % \includegraphics[width=0.4\textwidth]{images/logo_fmsp.png}
        \fbox{\parbox{0.4\textwidth}{\centering [LOGO DE L'APPLICATION]\\ \vspace{0.5cm} FMSP}}
    \end{figure}

    \vspace{2cm}
    {\Large Version 2.0 - Corrigée et détaillée\par}
    \vspace{0.5cm}
    {\large Basée sur l'analyse du code source\par}
    \vspace{1cm}
    {\large \today\par}

    \vfill
    {\large Faculté de Médecine et Sciences Pharmaceutiques\par}
\end{titlepage}

% Table des matières
\tableofcontents
\newpage

% Inclusion des sections
% Section 01 - Introduction
\section{Introduction}

\subsection{À propos de cette application}

Ce logiciel est une application de bureau développée spécifiquement pour la Faculté de Médecine et Sciences Pharmaceutiques (FMSP). Il permet de générer automatiquement :

\begin{itemize}[leftmargin=*]
    \item \textbf{Relevés de notes} détaillés pour les étudiants, avec calcul automatique des moyennes
    \item \textbf{Attestations de réussite} pour les étudiants ayant validé leurs semestres (moyenne $\geq$ 10/20)
    \item \textbf{Documents avec codes QR} pour faciliter la vérification et l'authenticité
\end{itemize}

L'application fonctionne entièrement \textbf{hors ligne} et vos données restent strictement confidentielles sur votre ordinateur.

\subsection{Fonctionnalités principales}

\subsubsection{Génération de documents}

\begin{description}[leftmargin=3cm, style=nextline]
    \item[Relevés de notes]
    \begin{itemize}
        \item Import de fichiers Excel avec données des étudiants
        \item Calcul automatique des moyennes par UE et moyenne générale
        \item Attribution automatique des mentions (Très Bien, Bien, Assez Bien, Passable, Ajourné)
        \item Gestion des crédits ECTS
        \item Support des sessions de rattrapage pour certains ECs
        \item QR code chiffré optionnel pour sécurité
        \item Génération par semestre ou multi-semestres
    \end{itemize}

    \item[Attestations de réussite]
    \begin{itemize}
        \item Validation automatique : seuls les étudiants avec moyenne $\geq$ 10/20 sont éligibles
        \item Statistiques en temps réel des étudiants admis/ajournés
        \item Thèmes entièrement personnalisables (couleurs, polices, mise en page)
        \item 6 préréglages de thèmes professionnels
        \item Support bilingue français/anglais
        \item QR code chiffré compact
    \end{itemize}

    \item[Codes QR sur documents]
    \begin{itemize}
        \item Ajout de QR codes sur PDFs existants (relevés, attestations, autres)
        \item QR code appliqué sur \textbf{TOUTES les pages} de chaque document
        \item Positionnement visuel précis (grille A4)
        \item Tailles personnalisables : petit, moyen, grand, ou taille personnalisée
        \item Niveaux de correction d'erreur ajustables (L, M, Q, H)
        \item Traitement en lot depuis Excel
    \end{itemize}
\end{description}

\subsubsection{Formats d'exportation}

L'application propose \textbf{3 formats d'exportation} pour tous les types de documents :

\begin{table}[h]
\centering
\begin{tabular}{|l|p{10cm}|}
\hline
\textbf{Format} & \textbf{Description} \\
\hline
\textbf{ZIP} & Archive compressée contenant tous les PDFs individuels. Format par défaut recommandé pour la distribution. \\
\hline
\textbf{Fichiers individuels} & Téléchargement séparé de chaque PDF (sans archive). Utile pour traiter un petit nombre de documents. \\
\hline
\textbf{PDF unique} & Tous les documents fusionnés en un seul fichier PDF. Idéal pour l'archivage centralisé. \\
\hline
\end{tabular}
\caption{Formats d'exportation disponibles}
\end{table}

\subsection{Les 9 menus de l'application}

L'interface de l'application est organisée en 9 menus accessibles depuis la barre latérale :

\begin{enumerate}[leftmargin=*]
    \item \menu{Générer les relevés} (\textit{Icône : FileSpreadsheet})
    \begin{itemize}
        \item Génération de relevés de notes à partir de fichiers Excel
        \item Composant : \texttt{ReleveGenerator}
    \end{itemize}

    \item \menu{Générer les attestations} (\textit{Icône : GraduationCap})
    \begin{itemize}
        \item Génération d'attestations de réussite
        \item Validation automatique des moyennes
        \item Composant : \texttt{AttestationGenerator}
    \end{itemize}

    \item \menu{Config des relevés} (\textit{Icône : Settings2})
    \begin{itemize}
        \item Configuration académique : semestres, UEs, ECs
        \item Gestion des coefficients et crédits
        \item Composant : \texttt{AcademicConfigManager}
    \end{itemize}

    \item \menu{Configurer les entêtes} (\textit{Icône : Cog})
    \begin{itemize}
        \item Paramètres de l'établissement (nom, logo, adresse)
        \item Configuration des informations affichées sur les documents
        \item Composant : \texttt{SettingsForm}
    \end{itemize}

    \item \menu{Historique des docs} (\textit{Icône : History})
    \begin{itemize}
        \item Consultation de l'historique des documents générés
        \item Recherche et filtrage
        \item Composant : \texttt{DocumentHistoryManager}
    \end{itemize}

    \item \menu{QR Codes sur PDF} (\textit{Icône : File})
    \begin{itemize}
        \item Ajout de QR codes sur PDFs existants
        \item Matching automatique avec données Excel
        \item Composant : \texttt{QrCodeOnPdf}
        \item \attention{Désactivé en mode démo}
    \end{itemize}

    \item \menu{Placement QR sur Documents} (\textit{Icône : QrCode})
    \begin{itemize}
        \item Placement avancé de QR codes avec mode manuel et Excel
        \item Personnalisation complète du QR code
        \item Composant : \texttt{QRCodeDocumentProcessor}
    \end{itemize}

    \item \menu{Aide} (\textit{Icône : HelpCircle})
    \begin{itemize}
        \item Documentation et support
        \item Composant : \texttt{HelpSupport}
    \end{itemize}

    \item \menu{Paramètres} (\textit{Icône : Settings})
    \begin{itemize}
        \item Configuration de l'application
        \item Gestion de la licence d'activation
        \item Composant : \texttt{SettingsPage}
    \end{itemize}
\end{enumerate}

\begin{figure}[h]
    \centering
    % \includegraphics[width=\textwidth]{images/interface_principale.png}
    \fbox{\parbox{0.9\textwidth}{\centering [CAPTURE D'ÉCRAN : INTERFACE PRINCIPALE AVEC SIDEBAR]\\ \vspace{0.3cm} Montrant les 9 menus dans la barre latérale gauche}}
    \caption{Interface principale de l'application}
\end{figure}

\subsection{Configuration requise}

\begin{tcolorbox}[colback=blue!5!white,colframe=blue!75!black,title=Configuration système minimale]
\begin{itemize}[leftmargin=*]
    \item \textbf{Système d'exploitation} : Windows 10/11, macOS 10.14+, ou Linux (Ubuntu 18.04+)
    \item \textbf{Processeur} : Intel Core i3 ou équivalent (2 GHz minimum)
    \item \textbf{Mémoire RAM} : 4 Go minimum (8 Go recommandé pour génération massive)
    \item \textbf{Espace disque} : 500 Mo pour l'application + espace pour documents générés
    \item \textbf{Résolution écran} : 1280x720 minimum (1920x1080 recommandé)
\end{itemize}
\end{tcolorbox}

\begin{tcolorbox}[colback=green!5!white,colframe=green!75!black,title=Configuration recommandée pour usage intensif]
\begin{itemize}[leftmargin=*]
    \item \textbf{Processeur} : Intel Core i5 ou supérieur (quad-core)
    \item \textbf{Mémoire RAM} : 8 Go ou plus
    \item \textbf{Disque dur} : SSD pour de meilleures performances
    \item \textbf{Générateur PDF} : Installation de Chrome ou Chromium (pour conversion HTML vers PDF)
\end{itemize}
\end{tcolorbox}

\subsection{Formats de fichiers supportés}

\subsubsection{Fichiers en entrée}

\begin{itemize}[leftmargin=*]
    \item \textbf{Données étudiants} : Excel (.xlsx, .xls) ou CSV
    \item \textbf{Documents PDF} : Tous fichiers PDF standards (pour ajout de QR codes)
    \item \textbf{Configurations} : Fichiers JSON (export/import de configurations)
    \item \textbf{Thèmes} : Fichiers JSON (export/import de thèmes)
\end{itemize}

\subsubsection{Fichiers générés en sortie}

\begin{itemize}[leftmargin=*]
    \item \textbf{Relevés de notes} : PDF (A4 portrait)
    \item \textbf{Attestations de réussite} : PDF (A4 portrait ou paysage selon thème)
    \item \textbf{Archives} : ZIP (compression optionnelle)
    \item \textbf{Historique} : Fichiers JSON de traçabilité
\end{itemize}

\subsection{Sécurité et confidentialité}

\begin{tcolorbox}[colback=yellow!10!white,colframe=orange!75!black,title=Protection des données]
\textbf{Garanties de confidentialité :}
\begin{itemize}[leftmargin=*]
    \item \textbf{Fonctionnement 100\% hors ligne} : Aucune connexion internet requise
    \item \textbf{Stockage local uniquement} : Toutes les données restent sur votre ordinateur
    \item \textbf{Aucun envoi de données} : Aucune information n'est transmise à l'extérieur
    \item \textbf{QR codes chiffrés} : Utilisation de AES-128-ECB avec clé basée sur matricule
    \item \textbf{Historique local} : Traçabilité stockée localement dans le navigateur
\end{itemize}

\textbf{Recommandations :}
\begin{itemize}[leftmargin=*]
    \item Effectuez des sauvegardes régulières de vos configurations
    \item Conservez les fichiers Excel sources en lieu sûr
    \item Exportez vos thèmes personnalisés pour éviter toute perte
    \item Sauvegardez l'historique des documents si nécessaire
\end{itemize}
\end{tcolorbox}

\subsection{Mode démo et licence}

L'application propose deux modes de fonctionnement :

\begin{description}[leftmargin=3cm]
    \item[Mode démo]
    \begin{itemize}
        \item Accès à la plupart des fonctionnalités
        \item Limitation : le menu \menu{QR Codes sur PDF} est désactivé
        \item Idéal pour tester l'application avant activation
    \end{itemize}

    \item[Mode activé]
    \begin{itemize}
        \item Accès complet à toutes les fonctionnalités
        \item Nécessite une clé de licence valide
        \item Configuration via le menu \menu{Paramètres}
    \end{itemize}
\end{description}

\newpage

% Section 02 - Installation et démarrage
\section{Installation et premier démarrage}

\subsection{Téléchargement de l'application}

\subsubsection{Où télécharger ?}

L'application est distribuée sous forme d'installateur pour chaque système d'exploitation :

\begin{itemize}[leftmargin=*]
    \item \textbf{Windows} : Fichier \texttt{FMSP-Releves-Setup-x.x.x.exe}
    \item \textbf{macOS} : Fichier \texttt{FMSP-Releves-x.x.x.dmg}
    \item \textbf{Linux} : Fichier \texttt{FMSP-Releves-x.x.x.AppImage} ou \texttt{.deb}
\end{itemize}

\astuce{Vérifiez toujours que vous téléchargez la dernière version depuis la source officielle.}

\subsection{Installation sur Windows}

\begin{enumerate}[leftmargin=*]
    \item Double-cliquez sur le fichier \texttt{FMSP-Releves-Setup-x.x.x.exe}
    \item Si Windows SmartScreen affiche un avertissement :
    \begin{itemize}
        \item Cliquez sur \bouton{Informations complémentaires}
        \item Puis sur \bouton{Exécuter quand même}
    \end{itemize}
    \item Suivez l'assistant d'installation :
    \begin{itemize}
        \item Acceptez la licence
        \item Choisissez le répertoire d'installation (défaut : \texttt{C:\textbackslash Program Files\textbackslash FMSP-Releves})
        \item Cochez "Créer un raccourci sur le bureau" si souhaité
        \item Cliquez sur \bouton{Installer}
    \end{itemize}
    \item Attendez la fin de l'installation (environ 30 secondes)
    \item Cliquez sur \bouton{Terminer} pour lancer l'application
\end{enumerate}

\begin{figure}[h]
    \centering
    % \includegraphics[width=0.7\textwidth]{images/installation_windows.png}
    \fbox{\parbox{0.7\textwidth}{\centering [CAPTURE D'ÉCRAN : ASSISTANT D'INSTALLATION WINDOWS]}}
    \caption{Assistant d'installation sous Windows}
\end{figure}

\subsection{Installation sur macOS}

\begin{enumerate}[leftmargin=*]
    \item Double-cliquez sur le fichier \texttt{FMSP-Releves-x.x.x.dmg}
    \item Une fenêtre s'ouvre montrant l'icône de l'application
    \item Glissez-déposez l'icône vers le dossier \texttt{Applications}
    \item Éjectez le volume DMG
    \item Ouvrez le dossier \texttt{Applications}
    \item Double-cliquez sur \texttt{FMSP-Releves}
    \item Si macOS affiche "Application non vérifiée" :
    \begin{itemize}
        \item Allez dans \menu{Préférences Système} > \menu{Sécurité et confidentialité}
        \item Cliquez sur \bouton{Ouvrir quand même}
    \end{itemize}
\end{enumerate}

\subsection{Installation sur Linux}

\subsubsection{Méthode 1 : AppImage (recommandé)}

\begin{enumerate}[leftmargin=*]
    \item Rendez le fichier exécutable :
    \begin{verbatim}
chmod +x FMSP-Releves-x.x.x.AppImage
    \end{verbatim}
    \item Lancez l'application :
    \begin{verbatim}
./FMSP-Releves-x.x.x.AppImage
    \end{verbatim}
\end{enumerate}

\subsubsection{Méthode 2 : Paquet Debian (.deb)}

\begin{enumerate}[leftmargin=*]
    \item Installez le paquet :
    \begin{verbatim}
sudo dpkg -i FMSP-Releves-x.x.x.deb
    \end{verbatim}
    \item Si des dépendances manquent :
    \begin{verbatim}
sudo apt-get install -f
    \end{verbatim}
    \item Lancez depuis le menu des applications ou en ligne de commande :
    \begin{verbatim}
fmsp-releves
    \end{verbatim}
\end{enumerate}

\subsection{Premier démarrage}

\subsubsection{Écran de bienvenue}

Au premier lancement, l'application affiche un écran de bienvenue vous guidant à travers les étapes initiales.

\begin{figure}[h]
    \centering
    % \includegraphics[width=0.8\textwidth]{images/ecran_bienvenue.png}
    \fbox{\parbox{0.8\textwidth}{\centering [CAPTURE D'ÉCRAN : ÉCRAN DE BIENVENUE]}}
    \caption{Écran de bienvenue au premier démarrage}
\end{figure}

\subsubsection{Configuration initiale de l'établissement}

Avant de générer vos premiers documents, configurez les informations de votre établissement :

\begin{enumerate}[leftmargin=*]
    \item Cliquez sur le menu \menu{Configurer les entêtes}
    \item Remplissez les informations suivantes :

    \begin{tcolorbox}[colback=blue!5!white,colframe=blue!75!black,title=Informations requises]
    \begin{itemize}[leftmargin=*]
        \item \champ{Nom de l'établissement} : Ex. "Faculté de Médecine et Sciences Pharmaceutiques"
        \item \champ{Adresse complète} : Rue, ville, code postal, pays
        \item \champ{Téléphone} : Numéro de contact
        \item \champ{Email} : Adresse email officielle
        \item \champ{Site web} : URL du site (optionnel)
        \item \champ{Logo} : Image PNG ou JPG (max 2 Mo, résolution recommandée : 800x200 px)
    \end{itemize}
    \end{tcolorbox}

    \item Cliquez sur \bouton{Enregistrer les paramètres}
\end{enumerate}

\begin{figure}[h]
    \centering
    % \includegraphics[width=\textwidth]{images/config_entetes.png}
    \fbox{\parbox{0.9\textwidth}{\centering [CAPTURE D'ÉCRAN : FORMULAIRE DE CONFIGURATION DES ENTÊTES]}}
    \caption{Configuration des entêtes de l'établissement}
\end{figure}

\subsection{Activation de la licence}

\subsubsection{Mode démo vs Mode activé}

\begin{table}[h]
\centering
\begin{tabular}{|l|c|c|}
\hline
\textbf{Fonctionnalité} & \textbf{Mode démo} & \textbf{Mode activé} \\
\hline
Générer les relevés & ✓ & ✓ \\
Générer les attestations & ✓ & ✓ \\
Config des relevés & ✓ & ✓ \\
Configurer les entêtes & ✓ & ✓ \\
Historique des docs & ✓ & ✓ \\
\textbf{QR Codes sur PDF} & \textbf{✗} & \textbf{✓} \\
Placement QR sur Documents & ✓ & ✓ \\
Aide & ✓ & ✓ \\
Paramètres & ✓ & ✓ \\
\hline
\end{tabular}
\caption{Comparaison mode démo vs mode activé}
\end{table}

\attention{Le menu "QR Codes sur PDF" est désactivé en mode démo. Pour l'activer, vous devez entrer une clé de licence valide.}

\subsubsection{Procédure d'activation}

\begin{enumerate}[leftmargin=*]
    \item Cliquez sur le menu \menu{Paramètres}
    \item Dans l'onglet \textbf{Licence}, vous verrez l'état actuel (Démo ou Activé)
    \item Entrez votre clé de licence dans le champ prévu
    \item Cliquez sur \bouton{Activer}
    \item Si la clé est valide, un message de confirmation s'affiche
    \item Redémarrez l'application pour appliquer l'activation
\end{enumerate}

\begin{figure}[h]
    \centering
    % \includegraphics[width=0.8\textwidth]{images/activation_licence.png}
    \fbox{\parbox{0.8\textwidth}{\centering [CAPTURE D'ÉCRAN : ÉCRAN D'ACTIVATION DE LA LICENCE]}}
    \caption{Interface d'activation de la licence}
\end{figure}

\subsection{Structure des dossiers de l'application}

Après installation, l'application crée automatiquement les dossiers suivants sur votre système :

\subsubsection{Windows}

\begin{verbatim}
C:\Users\[VotreNom]\AppData\Roaming\FMSP-Releves\
├── configs\              # Configurations de classes sauvegardées
├── themes\               # Thèmes personnalisés
├── history\              # Historique des documents générés
└── cache\                # Cache temporaire
\end{verbatim}

\subsubsection{macOS}

\begin{verbatim}
~/Library/Application Support/FMSP-Releves/
├── configs/
├── themes/
├── history/
└── cache/
\end{verbatim}

\subsubsection{Linux}

\begin{verbatim}
~/.config/FMSP-Releves/
├── configs/
├── themes/
├── history/
└── cache/
\end{verbatim}

\info{Ces dossiers contiennent vos données personnalisées. N'hésitez pas à en faire des sauvegardes régulières.}

\subsection{Vérification de l'installation}

Pour vérifier que l'application fonctionne correctement :

\begin{enumerate}[leftmargin=*]
    \item L'interface principale s'affiche correctement avec les 9 menus dans la barre latérale
    \item Le menu \menu{Aide} affiche la documentation
    \item Le menu \menu{Paramètres} montre l'état de la licence
    \item Vous pouvez accéder à tous les menus (sauf "QR Codes sur PDF" si en mode démo)
\end{enumerate}

\begin{tcolorbox}[colback=red!5!white,colframe=red!75!black,title=En cas de problème au démarrage]
Si l'application ne se lance pas correctement :
\begin{itemize}[leftmargin=*]
    \item Vérifiez que votre système répond à la configuration minimale
    \item Désinstallez et réinstallez l'application
    \item Vérifiez les logs d'erreur dans le dossier de l'application
    \item Consultez le menu \menu{Aide} > \menu{Rapporter un problème}
\end{itemize}
\end{tcolorbox}

\subsection{Désinstallation}

\subsubsection{Windows}

\begin{enumerate}[leftmargin=*]
    \item Panneau de configuration > Programmes et fonctionnalités
    \item Recherchez "FMSP-Releves"
    \item Clic droit > Désinstaller
    \item Suivez l'assistant de désinstallation
\end{enumerate}

\subsubsection{macOS}

\begin{enumerate}[leftmargin=*]
    \item Ouvrez le dossier Applications
    \item Glissez l'icône FMSP-Releves vers la Corbeille
    \item Videz la Corbeille
\end{enumerate}

\subsubsection{Linux}

Pour AppImage : supprimez simplement le fichier .AppImage

Pour paquet Debian :
\begin{verbatim}
sudo apt-get remove fmsp-releves
\end{verbatim}

\attention{La désinstallation ne supprime pas automatiquement vos données personnelles (configurations, thèmes, historique). Ces dossiers doivent être supprimés manuellement si souhaité.}

\newpage

% Section 03 - Génération de relevés de notes
\section{Génération de relevés de notes}

\subsection{Vue d'ensemble}

Le menu \menu{Générer les relevés} permet de créer automatiquement des relevés de notes personnalisés pour vos étudiants à partir d'un fichier Excel contenant les notes.

\textbf{Prérequis :}
\begin{itemize}[leftmargin=*]
    \item Avoir configuré au moins une classe dans \menu{Config des relevés}
    \item Disposer d'un fichier Excel avec les données des étudiants et leurs notes
    \item Avoir configuré les entêtes de l'établissement
\end{itemize}

\subsection{Préparation du fichier Excel}

\subsubsection{Structure obligatoire}

Votre fichier Excel doit contenir au minimum les colonnes suivantes :

\begin{tcolorbox}[colback=blue!5!white,colframe=blue!75!black,title=Colonnes obligatoires]
\begin{longtable}{|l|p{10cm}|}
\hline
\textbf{Colonne} & \textbf{Description} \\
\hline
\endfirsthead
\hline
\textbf{Colonne} & \textbf{Description} \\
\hline
\endhead
\hline
\endfoot

\texttt{MATRICULE} & Numéro matricule unique de l'étudiant \\
\hline
\texttt{NOM} & Nom de famille (sera affiché en MAJUSCULES) \\
\hline
\texttt{PRENOM} & Prénom(s) de l'étudiant \\
\hline
\texttt{SEMESTRE} & Code du semestre (ex: S1, S2, S3, S4, S5, S6) \\
\hline
\texttt{NIVEAU} & Niveau d'études (ex: L1, L2, L3, M1, M2) \\
\hline
\texttt{FILIERE} & Code ou intitulé de la filière \\
\hline
\end{longtable}
\end{tcolorbox}

\subsubsection{Colonnes de notes}

En plus des colonnes obligatoires, ajoutez une colonne pour chaque EC (Élément Constitutif) défini dans votre configuration de classe.

\textbf{Exemple :} Si votre configuration contient les ECs suivants :
\begin{itemize}
    \item Mathématiques
    \item Physique
    \item Chimie
    \item Biologie
\end{itemize}

Votre Excel devra contenir les colonnes :
\begin{verbatim}
MATRICULE | NOM | PRENOM | SEMESTRE | NIVEAU | FILIERE | Mathématiques |
Physique | Chimie | Biologie
\end{verbatim}

\begin{figure}[h]
    \centering
    % \includegraphics[width=\textwidth]{images/excel_structure_releves.png}
    \fbox{\parbox{0.95\textwidth}{\centering [CAPTURE D'ÉCRAN : EXEMPLE DE FICHIER EXCEL POUR RELEVÉS]\\ \vspace{0.3cm} Montrant les colonnes obligatoires + colonnes de notes}}
    \caption{Structure d'un fichier Excel pour génération de relevés}
\end{figure}

\subsubsection{Colonnes optionnelles}

Vous pouvez ajouter des colonnes supplémentaires qui ne seront pas utilisées par l'application :

\begin{itemize}[leftmargin=*]
    \item \texttt{DATE\_NAISSANCE} : Date de naissance
    \item \texttt{LIEU\_NAISSANCE} : Lieu de naissance
    \item \texttt{SEXE} : M/F
    \item \texttt{EMAIL} : Adresse email
    \item Toute autre colonne pour votre usage interne
\end{itemize}

\attention{Nommez vos colonnes de notes EXACTEMENT comme les noms d'ECs dans votre configuration de classe (sensible à la casse).}

\subsubsection{Format des notes}

\begin{itemize}[leftmargin=*]
    \item Les notes doivent être des nombres (format numérique Excel)
    \item Les notes peuvent être décimales : \texttt{15.5}, \texttt{12.75}
    \item Pour une note manquante, laissez la cellule vide ou mettez \texttt{ABS} (absent)
    \item Les notes doivent correspondre à la base définie dans la configuration (sur 10, 20, ou 100)
\end{itemize}

\subsection{Workflow complet de génération}

\subsubsection{Étape 1 : Accéder au menu}

\begin{enumerate}[leftmargin=*]
    \item Cliquez sur \menu{Générer les relevés} dans la barre latérale
    \item L'interface de génération s'affiche avec plusieurs onglets
\end{enumerate}

\subsubsection{Étape 2 : Charger le fichier Excel}

\begin{enumerate}[leftmargin=*]
    \item Dans l'onglet \textbf{Générateur}, cliquez sur \bouton{Charger fichier Excel}
    \item Sélectionnez votre fichier .xlsx ou .xls
    \item L'application valide automatiquement :
    \begin{itemize}
        \item Présence des colonnes obligatoires
        \item Format des données
        \item Nombre d'étudiants détectés
    \end{itemize}
    \item Un message de confirmation s'affiche avec le nombre de lignes importées
\end{enumerate}

\begin{figure}[h]
    \centering
    % \includegraphics[width=0.9\textwidth]{images/import_excel_releves.png}
    \fbox{\parbox{0.85\textwidth}{\centering [CAPTURE D'ÉCRAN : IMPORT EXCEL AVEC VALIDATION]\\ \vspace{0.3cm} Message : "✓ 150 étudiants importés avec succès"}}
    \caption{Import et validation du fichier Excel}
\end{figure}

\subsubsection{Étape 3 : Sélectionner la configuration de classe}

\begin{enumerate}[leftmargin=*]
    \item Dans le champ \champ{Sélectionner une configuration}, choisissez la classe correspondante
    \item Exemple : "Licence 1 Médecine - 2023-2024"
    \item La liste des semestres disponibles s'affiche automatiquement
\end{enumerate}

\info{Si vous n'avez pas encore créé de configuration, cliquez sur \bouton{Créer une nouvelle configuration} pour être redirigé vers le menu \menu{Config des relevés}.}

\subsubsection{Étape 4 : Choisir le(s) semestre(s)}

Vous avez deux options :

\paragraph{Option A : Un seul semestre}

\begin{enumerate}[leftmargin=*]
    \item Dans la liste déroulante \champ{Sélectionner un semestre}, choisissez le semestre souhaité
    \item Exemple : "Semestre 1 (S1)"
    \item Seuls les étudiants de ce semestre seront traités
\end{enumerate}

\paragraph{Option B : Tous les semestres (génération multi-semestres)}

\begin{enumerate}[leftmargin=*]
    \item Sélectionnez \textbf{"Tous les semestres"} dans la liste déroulante
    \item Une interface de sélection multi-semestres apparaît
    \item Cochez les semestres à inclure dans la génération
    \item Utilisez \bouton{Tout sélectionner} ou \bouton{Tout désélectionner} pour faciliter la sélection
\end{enumerate}

\begin{figure}[h]
    \centering
    % \includegraphics[width=\textwidth]{images/selection_multi_semestres.png}
    \fbox{\parbox{0.9\textwidth}{\centering [CAPTURE D'ÉCRAN : INTERFACE MULTI-SEMESTRES]\\ \vspace{0.3cm} Cases à cocher pour S1, S2, S3, S4, S5, S6\\ Boutons "Tout sélectionner" et "Tout désélectionner"}}
    \caption{Sélection de plusieurs semestres simultanément}
\end{figure}

\subsubsection{Étape 5 : Mapper les colonnes Excel aux ECs}

Si c'est la première fois que vous utilisez ce fichier Excel avec cette configuration, vous devez associer chaque EC à sa colonne Excel correspondante.

\begin{enumerate}[leftmargin=*]
    \item Cliquez sur l'onglet \textbf{Mapping des colonnes}
    \item Pour chaque EC listé, sélectionnez la colonne Excel correspondante dans la liste déroulante
    \item L'application propose automatiquement les colonnes dont le nom correspond
    \item Si un EC a une session de rattrapage :
    \begin{itemize}
        \item Activez le toggle \champ{Session}
        \item Mappez également la colonne de la session
    \end{itemize}
    \item Cliquez sur \bouton{Enregistrer le mapping}
\end{enumerate}

\begin{figure}[h]
    \centering
    % \includegraphics[width=\textwidth]{images/mapping_colonnes.png}
    \fbox{\parbox{0.9\textwidth}{\centering [CAPTURE D'ÉCRAN : INTERFACE DE MAPPING]\\ \vspace{0.3cm}
    Tableau avec colonnes : EC | Colonne Excel | Session | Colonne Session\\
    Mathématiques → [Dropdown: Mathématiques]\\
    Physique → [Dropdown: Physique Générale]}}
    \caption{Mapping des colonnes Excel vers les ECs de la configuration}
\end{figure}

\astuce{Le mapping est mémorisé ! La prochaine fois que vous chargerez un fichier Excel avec les mêmes noms de colonnes, le mapping se fera automatiquement.}

\subsubsection{Étape 6 : Configurer les options de génération}

Dans l'onglet \textbf{Options}, configurez les paramètres suivants :

\paragraph{Format d'exportation}

Choisissez parmi les 3 formats disponibles :

\begin{description}[leftmargin=3.5cm]
    \item[ZIP (défaut)] Archive compressée contenant tous les PDFs individuels
    \item[Fichiers individuels] Chaque relevé téléchargé séparément (sans archive)
    \item[PDF unique] Tous les relevés fusionnés en un seul fichier PDF
\end{description}

\paragraph{Compression}

\begin{itemize}[leftmargin=*]
    \item Activez \champ{Utiliser la compression} pour réduire la taille de l'archive ZIP
    \item Gain typique : 30-50\% de réduction de taille
    \item Légère augmentation du temps de génération
\end{itemize}

\paragraph{Format de nommage}

Choisissez comment les fichiers PDF seront nommés :

\begin{description}[leftmargin=3.5cm]
    \item[Détaillé (défaut)] \texttt{Releve\_NOM\_Prenom\_S1\_2023-2024.pdf}
    \item[Simple] \texttt{Releve\_MATRICULE.pdf}
\end{description}

\paragraph{QR Code chiffré (optionnel)}

\begin{itemize}[leftmargin=*]
    \item Activez \champ{Ajouter un QR code chiffré} pour intégrer un QR code sécurisé
    \item Le QR contient des informations chiffrées avec AES-128-ECB
    \item Clé de chiffrement basée sur le matricule de l'étudiant
    \item Permet la vérification d'authenticité du relevé
\end{itemize}

\begin{figure}[h]
    \centering
    % \includegraphics[width=0.8\textwidth]{images/options_generation_releves.png}
    \fbox{\parbox{0.75\textwidth}{\centering [CAPTURE D'ÉCRAN : ONGLET OPTIONS]\\ \vspace{0.3cm}
    Radio buttons : ○ ZIP ● Fichiers individuels ○ PDF unique\\
    ☑ Utiliser la compression\\
    Format nommage : [Dropdown: Détaillé]\\
    ☐ Ajouter un QR code chiffré}}
    \caption{Configuration des options de génération}
\end{figure}

\subsubsection{Étape 7 : Prévisualiser (optionnel)}

Avant de générer tous les relevés :

\begin{enumerate}[leftmargin=*]
    \item Cliquez sur l'onglet \textbf{Aperçu}
    \item Sélectionnez un étudiant dans la liste
    \item Cliquez sur \bouton{Générer l'aperçu}
    \item Un PDF de prévisualisation s'affiche dans le navigateur
    \item Vérifiez :
    \begin{itemize}
        \item Les informations de l'étudiant (nom, prénom, matricule)
        \item Les notes affichées pour chaque EC
        \item Le calcul des moyennes par UE
        \item La moyenne générale et la mention
        \item La mise en page et le thème
    \end{itemize}
    \item Si tout est correct, revenez à l'onglet \textbf{Générateur}
\end{enumerate}

\begin{figure}[h]
    \centering
    % \includegraphics[width=0.7\textwidth]{images/apercu_releve.png}
    \fbox{\parbox{0.65\textwidth}{\centering [CAPTURE D'ÉCRAN : APERÇU D'UN RELEVÉ PDF]\\ \vspace{0.3cm}
    Affichage d'un relevé complet avec notes, moyennes, mention}}
    \caption{Aperçu d'un relevé avant génération massive}
\end{figure}

\subsubsection{Étape 8 : Générer les relevés}

\begin{enumerate}[leftmargin=*]
    \item Retournez sur l'onglet \textbf{Générateur}
    \item Vérifiez que toutes les configurations sont correctes :
    \begin{itemize}
        \item Fichier Excel chargé ✓
        \item Configuration sélectionnée ✓
        \item Semestre(s) choisi(s) ✓
        \item Mapping des colonnes effectué ✓
        \item Options configurées ✓
    \end{itemize}
    \item Cliquez sur le bouton \bouton{Générer les relevés}
    \item Une barre de progression s'affiche montrant :
    \begin{itemize}
        \item Nombre de relevés générés / total
        \item Pourcentage de progression
        \item Temps estimé restant
    \end{itemize}
    \item Attendez la fin du traitement (temps variable selon le nombre d'étudiants)
\end{enumerate}

\begin{figure}[h]
    \centering
    % \includegraphics[width=0.8\textwidth]{images/progression_generation.png}
    \fbox{\parbox{0.75\textwidth}{\centering [CAPTURE D'ÉCRAN : BARRE DE PROGRESSION]\\ \vspace{0.3cm}
    Génération en cours : 47/150 relevés (31\%)\\
    ▓▓▓▓▓▓▓░░░░░░░░░░░░░░\\
    Temps estimé : 2 min 15 s}}
    \caption{Barre de progression pendant la génération}
\end{figure}

\subsubsection{Étape 9 : Téléchargement}

Une fois la génération terminée :

\begin{itemize}[leftmargin=*]
    \item \textbf{Format ZIP} : Une archive est automatiquement téléchargée
    \begin{itemize}
        \item Nom du fichier : \texttt{Releves\_[Date]\_[Heure].zip}
        \item Ouvrez l'archive et extrayez les PDFs
    \end{itemize}

    \item \textbf{Format Fichiers individuels} : Les fichiers sont téléchargés un par un
    \begin{itemize}
        \item Un délai de 100 ms entre chaque téléchargement évite de saturer le navigateur
        \item Les fichiers apparaissent dans votre dossier de téléchargements
    \end{itemize}

    \item \textbf{Format PDF unique} : Un seul fichier fusionné est téléchargé
    \begin{itemize}
        \item Nom : \texttt{Releves\_Fusionnes\_[Date].pdf}
        \item Contient tous les relevés dans l'ordre
    \end{itemize}
\end{itemize}

\subsection{Calcul automatique des moyennes et mentions}

\subsubsection{Calcul de la moyenne d'une UE}

Pour chaque Unité d'Enseignement (UE), la moyenne est calculée selon la formule :

\begin{tcolorbox}[colback=green!5!white,colframe=green!75!black,title=Formule de calcul]
$$\text{Moyenne UE} = \frac{\sum (\text{Note EC} \times \text{Poids EC})}{\sum \text{Poids EC}}$$

\textbf{Exemple :}

UE1 contient :
\begin{itemize}
    \item Mathématiques : note = 15/20, poids = 2
    \item Physique : note = 12/20, poids = 1.5
    \item Chimie : note = 14/20, poids = 2
\end{itemize}

$$\text{Moyenne UE1} = \frac{(15 \times 2) + (12 \times 1.5) + (14 \times 2)}{2 + 1.5 + 2} = \frac{30 + 18 + 28}{5.5} = \frac{76}{5.5} = 13.82$$
\end{tcolorbox}

\subsubsection{Calcul de la moyenne générale}

La moyenne générale du semestre est la moyenne pondérée des moyennes d'UE, pondérée par les crédits ECTS :

\begin{tcolorbox}[colback=green!5!white,colframe=green!75!black,title=Formule de calcul]
$$\text{Moyenne Générale} = \frac{\sum (\text{Moyenne UE} \times \text{Crédits UE})}{\sum \text{Crédits UE}}$$

\textbf{Exemple :}

\begin{itemize}
    \item UE1 : moyenne = 13.82, crédits = 30
    \item UE2 : moyenne = 11.5, crédits = 30
\end{itemize}

$$\text{Moyenne Générale} = \frac{(13.82 \times 30) + (11.5 \times 30)}{30 + 30} = \frac{414.6 + 345}{60} = \frac{759.6}{60} = 12.66$$
\end{tcolorbox}

\subsubsection{Attribution des mentions}

Les mentions sont attribuées automatiquement selon la moyenne générale :

\begin{table}[h]
\centering
\begin{tabular}{|l|c|l|}
\hline
\textbf{Mention} & \textbf{Moyenne} & \textbf{Couleur sur le relevé} \\
\hline
Très Bien & $\geq$ 16/20 & Vert foncé \\
\hline
Bien & 14 à 15.99 & Vert \\
\hline
Assez Bien & 12 à 13.99 & Bleu \\
\hline
Passable & 10 à 11.99 & Orange \\
\hline
Ajourné & < 10 & Rouge \\
\hline
\end{tabular}
\caption{Barème des mentions}
\end{table}

\subsection{Gestion des sessions de rattrapage}

Certains ECs peuvent avoir une session de rattrapage. Dans ce cas :

\begin{enumerate}[leftmargin=*]
    \item Lors de la configuration de l'EC dans \menu{Config des relevés}, cochez \champ{A une session}
    \item Dans le fichier Excel, créez deux colonnes :
    \begin{itemize}
        \item Une pour la note de session normale : \texttt{Mathematiques}
        \item Une pour la note de session de rattrapage : \texttt{Mathematiques\_Session}
    \end{itemize}
    \item Lors du mapping, associez les deux colonnes
    \item L'application prendra automatiquement la meilleure note entre les deux sessions
\end{enumerate}

\astuce{Si un étudiant n'a pas de note de rattrapage (cellule vide), seule la note de session normale sera utilisée.}

\subsection{Sélection d'étudiants spécifiques}

Vous pouvez générer des relevés pour certains étudiants uniquement :

\begin{enumerate}[leftmargin=*]
    \item Dans l'onglet \textbf{Sélection}, cochez \champ{Activer la sélection personnalisée}
    \item Trois options s'offrent à vous :

    \paragraph{Option 1 : Filtrer par niveau}
    \begin{itemize}
        \item Cochez les niveaux souhaités (L1, L2, L3, M1, M2...)
        \item Seuls les étudiants de ces niveaux seront traités
    \end{itemize}

    \paragraph{Option 2 : Filtrer par filière}
    \begin{itemize}
        \item Cochez les filières souhaitées
        \item Seuls les étudiants de ces filières seront traités
    \end{itemize}

    \paragraph{Option 3 : Sélection manuelle}
    \begin{itemize}
        \item Une liste complète des étudiants s'affiche
        \item Cochez individuellement les étudiants pour lesquels générer un relevé
        \item Utilisez la barre de recherche pour trouver rapidement un étudiant
    \end{itemize}

    \item Un compteur affiche le nombre d'étudiants sélectionnés
    \item Générez normalement avec \bouton{Générer les relevés}
\end{enumerate}

\begin{figure}[h]
    \centering
    % \includegraphics[width=\textwidth]{images/selection_etudiants.png}
    \fbox{\parbox{0.9\textwidth}{\centering [CAPTURE D'ÉCRAN : ONGLET SÉLECTION]\\ \vspace{0.3cm}
    ☑ Activer la sélection personnalisée\\
    Filtres : ☑ L1  ☐ L2  ☑ L3\\
    Liste avec cases à cocher par étudiant\\
    32 étudiants sélectionnés}}
    \caption{Sélection personnalisée d'étudiants}
\end{figure}

\subsection{Personnalisation du thème des relevés}

Vous pouvez personnaliser l'apparence des relevés de notes :

\begin{enumerate}[leftmargin=*]
    \item Dans le menu \menu{Config des relevés}, sélectionnez votre configuration
    \item Cliquez sur l'onglet \textbf{Thème}
    \item Configurez :
    \begin{itemize}
        \item \textbf{Couleurs} : Couleur primaire, secondaire, accent
        \item \textbf{Polices} : Famille de police, tailles
        \item \textbf{Mise en page} : Espacement, marges, disposition du contenu
    \end{itemize}
    \item Prévisualisez en temps réel les changements
    \item Enregistrez le thème
\end{enumerate}

Le thème sera automatiquement appliqué à tous les relevés générés avec cette configuration.

\newpage

% Section 04 - Génération d'attestations de réussite
\section{Génération d'attestations de réussite}

\subsection{Vue d'ensemble}

Le menu \menu{Générer les attestations} permet de créer des attestations officielles pour les étudiants ayant \textbf{validé} leur semestre (moyenne $\geq$ 10/20).

\attention{IMPORTANT : Seuls les étudiants avec une moyenne générale supérieure ou égale à 10/20 recevront une attestation. Les autres sont automatiquement exclus.}

\textbf{Caractéristiques principales :}
\begin{itemize}[leftmargin=*]
    \item Validation automatique des moyennes
    \item Thèmes entièrement personnalisables
    \item 6 préréglages professionnels
    \item Support bilingue français/anglais
    \item QR code chiffré compact
    \item Statistiques en temps réel
\end{itemize}

\subsection{Différences avec les relevés de notes}

\begin{table}[h]
\centering
\begin{tabular}{|l|p{5.5cm}|p{5.5cm}|}
\hline
\textbf{Critère} & \textbf{Relevés de notes} & \textbf{Attestations de réussite} \\
\hline
Public cible & Tous les étudiants & Étudiants admis uniquement ($\geq$ 10/20) \\
\hline
Contenu & Notes détaillées par EC, moyennes par UE & Résultat global (admis), pas de détail des notes \\
\hline
Personnalisation & Thème lié à la configuration & Thème indépendant et très personnalisable \\
\hline
Langues & Français uniquement & Bilingue FR/EN \\
\hline
QR Code & Chiffré AES-128-ECB & Chiffré compact \\
\hline
\end{tabular}
\caption{Comparaison relevés vs attestations}
\end{table}

\subsection{Préparation du fichier Excel}

\subsubsection{Colonnes obligatoires}

\begin{tcolorbox}[colback=blue!5!white,colframe=blue!75!black,title=Colonnes obligatoires]
\begin{longtable}{|l|p{9cm}|}
\hline
\textbf{Colonne} & \textbf{Description} \\
\hline
\endfirsthead
\hline
\textbf{Colonne} & \textbf{Description} \\
\hline
\endhead
\hline
\endfoot

\texttt{MATRICULE} & Numéro matricule unique \\
\hline
\texttt{NOM} & Nom de famille \\
\hline
\texttt{PRENOM} & Prénom(s) \\
\hline
\texttt{DATE\_NAISSANCE} & Date de naissance (format JJ/MM/AAAA ou DD/MM/YYYY) \\
\hline
\texttt{LIEU\_NAISSANCE} & Lieu de naissance \\
\hline
\texttt{ANNEE\_ACADEMIQUE} & Année académique (ex: 2023-2024) \\
\hline
\texttt{SEMESTRE} & Code du semestre (S1, S2, etc.) \\
\hline
\texttt{NIVEAU} & Niveau d'études (L1, L2, L3, M1, M2, etc.) \\
\hline
\texttt{DOMAINE} & Domaine d'études (ex: Sciences de la Santé) \\
\hline
\texttt{PARCOURS} & Parcours ou spécialité (ex: Médecine Générale) \\
\hline
\texttt{MOYENNE\_GENERALE} & Moyenne générale du semestre (sur 20) \\
\hline
\end{longtable}
\end{tcolorbox}

\subsubsection{Colonnes optionnelles pour support bilingue}

Si vous souhaitez générer des attestations en anglais, ajoutez les colonnes de traduction :

\begin{longtable}{|l|l|}
\hline
\textbf{Colonne française} & \textbf{Colonne anglaise} \\
\hline
\endfirsthead
\hline
\textbf{Colonne française} & \textbf{Colonne anglaise} \\
\hline
\endhead
\hline
\endfoot

\texttt{DOMAINE} & \texttt{DOMAINE\_EN} \\
\hline
\texttt{PARCOURS} & \texttt{PARCOURS\_EN} \\
\hline
\texttt{SPECIALITE} & \texttt{SPECIALITE\_EN} \\
\hline
\texttt{OPTION} & \texttt{OPTION\_EN} \\
\hline
\texttt{FINALITE} & \texttt{FINALITE\_EN} \\
\hline
\texttt{MENTION} & \texttt{MENTION\_EN} \\
\hline
\end{longtable}

\begin{figure}[h]
    \centering
    % \includegraphics[width=\textwidth]{images/excel_attestations.png}
    \fbox{\parbox{0.95\textwidth}{\centering [CAPTURE D'ÉCRAN : EXEMPLE FICHIER EXCEL POUR ATTESTATIONS]\\ \vspace{0.3cm}
    Montrant toutes les colonnes obligatoires + colonnes bilingues}}
    \caption{Structure du fichier Excel pour attestations de réussite}
\end{figure}

\subsection{Workflow complet de génération}

\subsubsection{Étape 1 : Accéder au menu}

\begin{enumerate}[leftmargin=*]
    \item Cliquez sur \menu{Générer les attestations} dans la barre latérale
    \item L'interface s'affiche avec 6 onglets :
    \begin{enumerate}
        \item \textbf{Générateur} : Import Excel et génération
        \item \textbf{Sélection} : Filtrage des étudiants éligibles
        \item \textbf{Préréglages} : Thèmes prédéfinis
        \item \textbf{Thème} : Personnalisation basique
        \item \textbf{Style Avancé} : Personnalisation typographique
        \item \textbf{Paramètres} : Options d'export et QR code
    \end{enumerate}
\end{enumerate}

\subsubsection{Étape 2 : Charger le fichier Excel}

\begin{enumerate}[leftmargin=*]
    \item Dans l'onglet \textbf{Générateur}, cliquez sur \bouton{Charger fichier Excel}
    \item Sélectionnez votre fichier contenant les données des étudiants
    \item L'application valide automatiquement :
    \begin{itemize}
        \item Présence des colonnes obligatoires
        \item Format de la colonne \texttt{MOYENNE\_GENERALE}
        \item Format des dates
    \end{itemize}
    \item Un message de confirmation s'affiche
\end{enumerate}

\subsubsection{Étape 3 : Analyse automatique de l'éligibilité}

Dès le chargement, l'application analyse automatiquement les moyennes :

\begin{tcolorbox}[colback=green!5!white,colframe=green!75!black,title=Statistiques affichées]
\textbf{Exemple d'affichage :}

\begin{itemize}[leftmargin=*]
    \item \textcolor{green}{\textbf{✓ 127 étudiants ADMIS}} (moyenne $\geq$ 10/20)
    \item \textcolor{red}{\textbf{✗ 23 étudiants AJOURNÉS}} (moyenne < 10/20)
    \item \textcolor{blue}{\textbf{Total : 150 étudiants}}
\end{itemize}

\textbf{Taux de réussite : 84.67\%}
\end{tcolorbox}

\begin{figure}[h]
    \centering
    % \includegraphics[width=0.7\textwidth]{images/stats_eligibilite.png}
    \fbox{\parbox{0.65\textwidth}{\centering [CAPTURE D'ÉCRAN : STATISTIQUES D'ÉLIGIBILITÉ]\\ \vspace{0.3cm}
    Graphique circulaire ou barres montrant admis/ajournés\\
    Pourcentages et nombres absolus}}
    \caption{Statistiques en temps réel des étudiants éligibles}
\end{figure}

\attention{Les étudiants ajournés n'apparaîtront pas dans la génération finale, même s'ils sont présents dans l'Excel.}

\subsubsection{Étape 4 : Choisir ou personnaliser un thème}

Vous avez deux options pour le design de vos attestations :

\paragraph{Option A : Utiliser un préréglage (rapide)}

\begin{enumerate}[leftmargin=*]
    \item Accédez à l'onglet \textbf{Préréglages}
    \item L'application propose 6 thèmes professionnels :
    \begin{itemize}
        \item \textbf{Classique Élégant} : Design sobre avec bleu marine et or
        \item \textbf{Moderne Minimaliste} : Design épuré avec gris et vert
        \item \textbf{Professionnel Académique} : Style traditionnel universitaire
        \item \textbf{Dynamique Coloré} : Couleurs vives et modernes
        \item \textbf{Excellence Prestige} : Design luxueux avec bordeaux et or
        \item \textbf{Nature Zen} : Tons verts et beiges apaisants
    \end{itemize}
    \item Cliquez sur \bouton{Aperçu} pour voir un exemple
    \item Cliquez sur \bouton{Appliquer ce thème} pour l'utiliser
\end{enumerate}

\begin{figure}[h]
    \centering
    % \includegraphics[width=\textwidth]{images/prereglages_themes.png}
    \fbox{\parbox{0.9\textwidth}{\centering [CAPTURE D'ÉCRAN : GALERIE DES PRÉRÉGLAGES]\\ \vspace{0.3cm}
    6 cartes affichant les aperçus des thèmes\\
    Boutons "Aperçu" et "Appliquer" sous chaque thème}}
    \caption{Galerie des 6 préréglages de thèmes disponibles}
\end{figure}

\paragraph{Option B : Personnaliser entièrement (avancé)}

\begin{enumerate}[leftmargin=*]
    \item Accédez à l'onglet \textbf{Thème}
    \item Configurez les éléments suivants :

    \subparagraph{Couleurs}
    \begin{itemize}
        \item \champ{Couleur primaire} : Couleur principale du design
        \item \champ{Couleur secondaire} : Couleur d'accentuation
        \item \champ{Couleur d'arrière-plan} : Fond de l'attestation
        \item \champ{Couleur du texte} : Texte principal
    \end{itemize}

    \subparagraph{En-tête}
    \begin{itemize}
        \item \champ{Afficher le logo} : Oui/Non
        \item \champ{Position du logo} : Gauche / Centre / Droite
        \item \champ{Taille du logo} : Petit / Moyen / Grand
        \item \champ{Afficher le nom de l'établissement} : Oui/Non
    \end{itemize}

    \subparagraph{Corps du document}
    \begin{itemize}
        \item \champ{Alignement du texte} : Gauche / Centre / Justifié
        \item \champ{Espacement entre lignes} : 1.0 / 1.5 / 2.0
        \item \champ{Afficher la mention} : Oui/Non
        \item \champ{Afficher la moyenne} : Oui/Non
    \end{itemize}

    \subparagraph{Pied de page}
    \begin{itemize}
        \item \champ{Signature} : Charger une image de signature
        \item \champ{Nom du signataire}
        \item \champ{Fonction du signataire}
        \item \champ{Date de délivrance} : Date du jour / Personnalisée
    \end{itemize}

    \item Prévisualisez en temps réel dans le panneau de droite
    \item Enregistrez votre thème personnalisé
\end{enumerate}

\begin{figure}[h]
    \centering
    % \includegraphics[width=\textwidth]{images/personnalisation_theme.png}
    \fbox{\parbox{0.9\textwidth}{\centering [CAPTURE D'ÉCRAN : INTERFACE DE PERSONNALISATION]\\ \vspace{0.3cm}
    Panneau gauche : Formulaire avec tous les réglages\\
    Panneau droit : Aperçu en temps réel de l'attestation}}
    \caption{Interface de personnalisation complète du thème}
\end{figure}

\subsubsection{Étape 5 : Configuration typographique avancée (optionnel)}

Pour un contrôle total sur la typographie :

\begin{enumerate}[leftmargin=*]
    \item Accédez à l'onglet \textbf{Style Avancé}
    \item Configurez :

    \paragraph{Polices de caractères}
    \begin{itemize}
        \item \champ{Famille de police pour titres} : Serif / Sans-serif / Monospace / Police personnalisée
        \item \champ{Famille de police pour le corps} : Idem
        \item \champ{Taille du titre principal} : 24pt à 48pt
        \item \champ{Taille des sous-titres} : 14pt à 24pt
        \item \champ{Taille du corps de texte} : 10pt à 16pt
    \end{itemize}

    \paragraph{Effets de texte}
    \begin{itemize}
        \item \champ{Style du titre} : Normal / Gras / Italique / Gras Italique
        \item \champ{Lettrage} : Normal / Petites capitales / Majuscules
        \item \champ{Ombrage du texte} : Oui/Non + couleur et décalage
    \end{itemize}

    \paragraph{Bordures et cadres}
    \begin{itemize}
        \item \champ{Bordure du document} : Oui/Non
        \item \champ{Style de bordure} : Simple / Double / Décorative
        \item \champ{Épaisseur} : 1px à 10px
        \item \champ{Couleur de bordure}
        \item \champ{Arrondi des coins} : 0px (carré) à 20px (très arrondi)
    \end{itemize}

    \item Les modifications sont visibles en temps réel dans l'aperçu
\end{enumerate}

\subsubsection{Étape 6 : Configurer les paramètres d'export}

\begin{enumerate}[leftmargin=*]
    \item Accédez à l'onglet \textbf{Paramètres}
    \item Configurez les options suivantes :

    \paragraph{Format d'exportation}
    \begin{description}[leftmargin=3.5cm]
        \item[ZIP (défaut)] Archive contenant tous les PDFs
        \item[Fichiers individuels] Téléchargement séparé de chaque attestation
        \item[PDF unique] Toutes les attestations fusionnées en un seul fichier
    \end{description}

    \paragraph{Compression}
    \begin{itemize}
        \item Activez \champ{Utiliser la compression} pour réduire la taille du ZIP
    \end{itemize}

    \paragraph{QR Code chiffré compact}
    \begin{itemize}
        \item Activez \champ{Ajouter un QR code chiffré}
        \item Le QR code contient :
        \begin{itemize}
            \item Matricule de l'étudiant
            \item Nom et prénom
            \item Semestre et année académique
            \item Hash de vérification
        \end{itemize}
        \item Chiffrement : AES-128-ECB
        \item Position : Coin inférieur droit par défaut
    \end{itemize}

    \paragraph{Langue}
    \begin{itemize}
        \item \champ{Langue de génération} : Français / Anglais
        \item Si Anglais sélectionné, les colonnes \texttt{*\_EN} de l'Excel seront utilisées
    \end{itemize}
\end{enumerate}

\begin{figure}[h]
    \centering
    % \includegraphics[width=0.8\textwidth]{images/parametres_attestations.png}
    \fbox{\parbox{0.75\textwidth}{\centering [CAPTURE D'ÉCRAN : ONGLET PARAMÈTRES]\\ \vspace{0.3cm}
    Format : ● ZIP  ○ Fichiers individuels  ○ PDF unique\\
    ☑ Utiliser la compression\\
    ☑ Ajouter un QR code chiffré\\
    Langue : [Dropdown: Français]}}
    \caption{Configuration des paramètres d'exportation}
\end{figure}

\subsubsection{Étape 7 : Filtrer les étudiants (optionnel)}

Par défaut, des attestations seront générées pour TOUS les étudiants admis. Pour restreindre :

\begin{enumerate}[leftmargin=*]
    \item Accédez à l'onglet \textbf{Sélection}
    \item Vous voyez la liste complète des étudiants éligibles (moyenne $\geq$ 10)
    \item Trois options de filtrage :

    \paragraph{Filtrer par niveau}
    \begin{itemize}
        \item Cochez les niveaux souhaités : L1, L2, L3, M1, M2, etc.
    \end{itemize}

    \paragraph{Filtrer par semestre}
    \begin{itemize}
        \item Cochez les semestres souhaités : S1, S2, S3, S4, S5, S6
    \end{itemize}

    \paragraph{Sélection manuelle}
    \begin{itemize}
        \item Cochez individuellement les étudiants
        \item Utilisez la barre de recherche pour trouver rapidement un étudiant
    \end{itemize}

    \item Un compteur affiche : \texttt{45/127 étudiants sélectionnés}
\end{enumerate}

\begin{figure}[h]
    \centering
    % \includegraphics[width=\textwidth]{images/selection_attestations.png}
    \fbox{\parbox{0.9\textwidth}{\centering [CAPTURE D'ÉCRAN : ONGLET SÉLECTION]\\ \vspace{0.3cm}
    Liste des étudiants admis avec cases à cocher\\
    Filtres par niveau et semestre en haut\\
    Barre de recherche\\
    Compteur "45/127 sélectionnés"}}
    \caption{Filtrage et sélection des étudiants pour attestations}
\end{figure}

\subsubsection{Étape 8 : Générer les attestations}

\begin{enumerate}[leftmargin=*]
    \item Retournez sur l'onglet \textbf{Générateur}
    \item Vérifiez le récapitulatif affiché :
    \begin{itemize}
        \item Nombre d'attestations à générer
        \item Thème sélectionné
        \item Format d'exportation
        \item Options activées (compression, QR code, etc.)
    \end{itemize}
    \item Cliquez sur \bouton{Générer les attestations}
    \item Une barre de progression s'affiche
    \item Attendez la fin du traitement
\end{enumerate}

\begin{figure}[h]
    \centering
    % \includegraphics[width=0.8\textwidth]{images/generation_attestations.png}
    \fbox{\parbox{0.75\textwidth}{\centering [CAPTURE D'ÉCRAN : GÉNÉRATION EN COURS]\\ \vspace{0.3cm}
    Génération de 127 attestations...\\
    ▓▓▓▓▓▓▓▓▓░░░░░░  63/127 (49\%)\\
    Temps restant estimé : 1 min 30 s}}
    \caption{Progression de la génération des attestations}
\end{figure}

\subsubsection{Étape 9 : Téléchargement et vérification}

Une fois la génération terminée :

\begin{itemize}[leftmargin=*]
    \item Un message de succès s'affiche : \texttt{✓ 127 attestations générées avec succès}
    \item Le fichier ou les fichiers sont téléchargés automatiquement
    \item Ouvrez quelques attestations pour vérifier :
    \begin{itemize}
        \item Les informations de l'étudiant sont correctes
        \item Le thème est appliqué correctement
        \item Le QR code est présent (si activé)
        \item La mise en page est satisfaisante
    \end{itemize}
\end{itemize}

\subsection{Export et import de thèmes personnalisés}

Pour réutiliser vos thèmes ou les partager :

\subsubsection{Exporter un thème}

\begin{enumerate}[leftmargin=*]
    \item Dans l'onglet \textbf{Thème}, configurez votre thème personnalisé
    \item Cliquez sur \bouton{Exporter le thème}
    \item Donnez un nom au thème : "Mon Thème Personnalisé"
    \item Un fichier JSON est téléchargé : \texttt{theme\_mon\_theme\_personnalise.json}
    \item Conservez ce fichier en lieu sûr
\end{enumerate}

\subsubsection{Importer un thème}

\begin{enumerate}[leftmargin=*]
    \item Dans l'onglet \textbf{Thème}, cliquez sur \bouton{Importer un thème}
    \item Sélectionnez le fichier JSON du thème
    \item Le thème est chargé et appliqué immédiatement
    \item Vous pouvez le modifier si nécessaire
\end{enumerate}

\astuce{Les thèmes exportés peuvent être partagés avec d'autres utilisateurs de l'application. Idéal pour standardiser l'apparence des attestations dans plusieurs services.}

\subsection{Support bilingue français/anglais}

\subsubsection{Activer le mode anglais}

\begin{enumerate}[leftmargin=*]
    \item Assurez-vous que votre fichier Excel contient les colonnes de traduction (\texttt{*\_EN})
    \item Dans l'onglet \textbf{Paramètres}, sélectionnez \champ{Langue : Anglais}
    \item Générez normalement
    \item Les attestations seront en anglais avec :
    \begin{itemize}
        \item Titre : "Certificate of Achievement" au lieu de "Attestation de Réussite"
        \item Contenu traduit automatiquement
        \item Valeurs des colonnes \texttt{*\_EN} utilisées
    \end{itemize}
\end{enumerate}

\subsubsection{Exemple de traduction}

\begin{table}[h]
\centering
\begin{tabular}{|l|l|}
\hline
\textbf{Français} & \textbf{Anglais (avec colonnes \_EN)} \\
\hline
Domaine: Sciences de la Santé & Field: Health Sciences \\
\hline
Parcours: Médecine Générale & Program: General Medicine \\
\hline
Spécialité: Cardiologie & Specialty: Cardiology \\
\hline
Mention: Très Bien & Grade: Excellent \\
\hline
\end{tabular}
\caption{Exemple de traduction automatique}
\end{table}

\newpage

% Section 05 - QR Code sur PDF
\section{Ajout de codes QR sur PDFs existants}

\subsection{Vue d'ensemble}

Le menu \menu{QR Codes sur PDF} permet d'ajouter des codes QR sur des documents PDF déjà générés (relevés, attestations, ou autres documents).

\attention{Cette fonctionnalité est DÉSACTIVÉE en mode démo. Vous devez activer l'application avec une licence valide pour y accéder.}

\textbf{Caractéristiques principales :}
\begin{itemize}[leftmargin=*]
    \item QR code ajouté sur \textbf{TOUTES les pages} de chaque PDF
    \item Matching automatique entre PDFs et données Excel
    \item Positionnement visuel précis avec grille A4
    \item Support des orientations portrait et paysage
    \item 3 types de documents supportés : relevé, attestation, diplôme
    \item 3 formats d'export : ZIP, Fichiers individuels, PDF unique
\end{itemize}

\subsection{Quand utiliser cette fonctionnalité ?}

\textbf{Cas d'usage typiques :}

\begin{enumerate}[leftmargin=*]
    \item Vous avez généré des relevés SANS QR code et souhaitez les ajouter après coup
    \item Vous voulez ajouter des QR codes sur des documents PDF externes (non générés par l'application)
    \item Vous devez mettre à jour les QR codes de documents existants
    \item Vous voulez ajouter des informations de vérification sur des documents scannés
\end{enumerate}

\info{Si vous générez des relevés ou attestations depuis l'application, vous pouvez directement activer l'option "QR code chiffré" dans les paramètres de génération. Cette fonctionnalité est pour ajouter des QR APRÈS la génération.}

\subsection{Workflow complet}

\subsubsection{Étape 1 : Accéder au menu}

\begin{enumerate}[leftmargin=*]
    \item Assurez-vous d'avoir une licence active (sinon le menu est grisé)
    \item Cliquez sur \menu{QR Codes sur PDF} dans la barre latérale
    \item L'interface s'affiche avec plusieurs sections
\end{enumerate}

\subsubsection{Étape 2 : Charger les documents PDF}

\begin{enumerate}[leftmargin=*]
    \item Cliquez sur \bouton{Charger des PDFs}
    \item Sélectionnez un ou plusieurs fichiers PDF
    \item L'application analyse chaque PDF :
    \begin{itemize}
        \item Nombre de pages détecté
        \item Taille du fichier
        \item Orientation (portrait/paysage)
    \end{itemize}
    \item Les PDFs apparaissent dans une liste
\end{enumerate}

\begin{figure}[h]
    \centering
    % \includegraphics[width=\textwidth]{images/liste_pdfs_qr.png}
    \fbox{\parbox{0.9\textwidth}{\centering [CAPTURE D'ÉCRAN : LISTE DES PDFs CHARGÉS]\\ \vspace{0.3cm}
    Tableau : Nom fichier | Pages | Taille | Statut | Actions\\
    Releve\_DUPONT\_Jean.pdf | 1 | 245 Ko | ✓ Prêt | [Supprimer]\\
    Releve\_MARTIN\_Marie.pdf | 2 | 387 Ko | ✓ Prêt | [Supprimer]}}
    \caption{Liste des PDFs chargés pour traitement}
\end{figure}

\subsubsection{Étape 3 : Charger le fichier Excel avec données}

\begin{enumerate}[leftmargin=*]
    \item Cliquez sur \bouton{Charger fichier Excel}
    \item Sélectionnez le fichier contenant les données des étudiants
    \item L'application lit les colonnes disponibles
\end{enumerate}

\textbf{Colonnes minimales requises :}
\begin{itemize}[leftmargin=*]
    \item \texttt{MATRICULE} : Pour le matching avec les noms de fichiers PDF
    \item \texttt{NOM} : Nom de l'étudiant
    \item \texttt{PRENOM} : Prénom de l'étudiant
    \item Toutes autres colonnes que vous souhaitez inclure dans le QR code
\end{itemize}

\subsubsection{Étape 4 : Vérifier le matching automatique}

L'application tente automatiquement d'associer chaque PDF à une ligne de l'Excel :

\begin{enumerate}[leftmargin=*]
    \item Le matching se fait par recherche du matricule dans le nom du fichier PDF
    \item Exemple : \texttt{Releve\_2023001\_DUPONT.pdf} → Ligne avec \texttt{MATRICULE = 2023001}
    \item Les correspondances trouvées sont marquées \textcolor{green}{\textbf{✓ Matched}}
    \item Les correspondances manquantes sont marquées \textcolor{red}{\textbf{✗ Non matché}}
\end{enumerate}

\begin{figure}[h]
    \centering
    % \includegraphics[width=\textwidth]{images/matching_pdfs.png}
    \fbox{\parbox}{0.9\textwidth}{\centering [CAPTURE D'ÉCRAN : RÉSULTAT DU MATCHING]\\ \vspace{0.3cm}
    Tableau avec statut de matching pour chaque PDF\\
    ✓ Vert pour les matchés\\
    ✗ Rouge pour les non matchés avec bouton "Associer manuellement"}}
    \caption{Résultat du matching automatique PDFs-Excel}
\end{figure}

\subsubsection{Étape 5 : Associer manuellement les PDFs non matchés (si nécessaire)}

Pour les PDFs non automatiquement matchés :

\begin{enumerate}[leftmargin=*]
    \item Cliquez sur \bouton{Associer manuellement} à côté du PDF
    \item Une fenêtre s'ouvre avec la liste de tous les étudiants de l'Excel
    \item Utilisez la barre de recherche pour trouver l'étudiant correct
    \item Cliquez sur \bouton{Associer} à côté de l'étudiant
    \item Le PDF passe au statut \textcolor{green}{✓ Matched}
\end{enumerate}

\subsubsection{Étape 6 : Sélectionner le type de document}

\begin{enumerate}[leftmargin=*]
    \item Dans le champ \champ{Type de document}, sélectionnez :
    \begin{description}[leftmargin=3cm]
        \item[Relevé] Pour des relevés de notes (orientation portrait A4)
        \item[Attestation] Pour des attestations (orientation portrait A4)
        \item[Diplôme] Pour des diplômes (orientation paysage A4)
    \end{description}
    \item Ce choix détermine l'orientation par défaut de la grille de positionnement
\end{enumerate}

\subsubsection{Étape 7 : Positionner le QR code}

\begin{enumerate}[leftmargin=*]
    \item Une grille de positionnement A4 s'affiche
    \item Cliquez directement sur la grille pour placer le QR code
    \item Ou utilisez les positions prédéfinies :
    \begin{itemize}
        \item \bouton{Haut Gauche}
        \item \bouton{Haut Droite}
        \item \bouton{Centre}
        \item \bouton{Bas Gauche}
        \item \bouton{Bas Droite}
    \end{itemize}
    \item Les coordonnées X et Y (en pourcentage) sont affichées
    \item Un aperçu visuel montre l'emplacement du QR sur la page
\end{enumerate}

\begin{figure}[h]
    \centering
    % \includegraphics[width=0.7\textwidth]{images/grille_positionnement_qr.png}
    \fbox{\parbox{0.65\textwidth}{\centering [CAPTURE D'ÉCRAN : GRILLE DE POSITIONNEMENT]\\ \vspace{0.3cm}
    Grille A4 cliquable avec aperçu du QR code\\
    Boutons positions prédéfinies\\
    Coordonnées affichées : X: 85\%, Y: 15\%}}
    \caption{Grille de positionnement visuel du QR code}
\end{figure}

\astuce{Pour un placement standard sur les relevés, utilisez la position "Haut Droite" (85\%, 15\%). Le QR code sera visible sans masquer d'informations importantes.}

\subsubsection{Étape 8 : Configurer le contenu du QR code}

\begin{enumerate}[leftmargin=*]
    \item Sélectionnez les colonnes Excel à inclure dans le QR code
    \item Recommandations :
    \begin{itemize}
        \item Minimum : MATRICULE, NOM, PRENOM
        \item Optionnel : SEMESTRE, NIVEAU, DATE\_NAISSANCE
        \item Évitez : Trop de données (QR code trop dense)
    \end{itemize}
    \item Cochez les colonnes souhaitées dans la liste
    \item Un aperçu du contenu QR s'affiche :
\end{enumerate}

\begin{verbatim}
Contenu du QR code:
MATRICULE: 2023001
NOM: DUPONT
PRENOM: Jean
SEMESTRE: S1
NIVEAU: L1
\end{verbatim}

\subsubsection{Étape 9 : Configurer les paramètres du QR code}

\begin{description}[leftmargin=3.5cm]
    \item[Taille] Choisissez parmi :
    \begin{itemize}
        \item Petit (60x60 pixels)
        \item Moyen (80x80 pixels) - Recommandé
        \item Grand (100x100 pixels)
        \item Personnalisée (saisissez une taille en pixels)
    \end{itemize}

    \item[Correction d'erreur] Niveau de récupération si le QR est endommagé :
    \begin{itemize}
        \item L (Low) : 7\% de récupération - QR plus simple
        \item M (Medium) : 15\% de récupération - \textbf{Recommandé}
        \item Q (Quartile) : 25\% de récupération
        \item H (High) : 30\% de récupération - QR plus dense
    \end{itemize}
\end{description}

\begin{figure}[h]
    \centering
    % \includegraphics[width=0.8\textwidth]{images/config_qr_params.png}
    \fbox{\parbox{0.75\textwidth}{\centering [CAPTURE D'ÉCRAN : PARAMÈTRES DU QR CODE]\\ \vspace{0.3cm}
    Taille : ○ Petit  ● Moyen  ○ Grand  ○ Personnalisée [ ]px\\
    Correction d'erreur : ○ L  ● M  ○ Q  ○ H\\
    Colonnes : ☑ MATRICULE  ☑ NOM  ☑ PRENOM  ☑ SEMESTRE}}
    \caption{Configuration des paramètres du QR code}
\end{figure}

\subsubsection{Étape 10 : Choisir le format d'exportation}

\begin{description}[leftmargin=3.5cm]
    \item[ZIP (défaut)] Archive contenant tous les PDFs avec QR codes
    \item[Fichiers individuels] Chaque PDF téléchargé séparément
    \item[PDF unique] Tous les PDFs fusionnés en un seul fichier
\end{description}

\subsubsection{Étape 11 : Traiter les documents}

\begin{enumerate}[leftmargin=*]
    \item Vérifiez le récapitulatif :
    \begin{itemize}
        \item Nombre de PDFs à traiter
        \item Nombre de PDFs matchés
        \item Position du QR code
        \item Colonnes incluses
        \item Format d'export
    \end{itemize}
    \item Cliquez sur \bouton{Traiter tous les documents}
    \item Une barre de progression s'affiche :
    \begin{itemize}
        \item Document en cours de traitement
        \item Numéro de page en cours (pour les PDFs multi-pages)
        \item Pourcentage global de progression
    \end{itemize}
\end{enumerate}

\begin{figure}[h]
    \centering
    % \includegraphics[width=0.8\textwidth]{images/traitement_qr_pdf.png}
    \fbox{\parbox{0.75\textwidth}{\centering [CAPTURE D'ÉCRAN : TRAITEMENT EN COURS]\\ \vspace{0.3cm}
    Traitement de Releve\_MARTIN\_Marie.pdf\\
    Page 2/2 en cours...\\
    ▓▓▓▓▓▓▓░░░░░░░░  23/50 PDFs (46\%)}}
    \caption{Barre de progression du traitement}
\end{figure}

\info{Le QR code est ajouté sur TOUTES les pages de chaque PDF. Pour un relevé de 2 pages, le même QR code apparaîtra sur les 2 pages.}

\subsubsection{Étape 12 : Téléchargement et vérification}

\begin{enumerate}[leftmargin=*]
    \item Une fois terminé, un message de confirmation s'affiche :
    \begin{itemize}
        \item \textcolor{green}{✓ 50 documents traités avec succès}
        \item \textcolor{red}{✗ 0 erreur}
    \end{itemize}
    \item Le fichier ou les fichiers sont téléchargés automatiquement
    \item Ouvrez quelques PDFs pour vérifier :
    \begin{itemize}
        \item Le QR code est présent sur toutes les pages
        \item La position est correcte
        \item Le QR code est lisible (testez avec un scanner)
        \item Le QR ne masque pas d'informations importantes
    \end{itemize}
\end{enumerate}

\subsection{Traitement de PDFs multi-pages}

\textbf{Particularité importante :} Contrairement à la plupart des outils, cette application ajoute le QR code sur **TOUTES les pages** de chaque document.

\textbf{Exemple :}
\begin{itemize}[leftmargin=*]
    \item Vous avez un relevé de 3 pages
    \item Le QR code sera ajouté identiquement sur la page 1, page 2 ET page 3
    \item Chaque page du document peut être vérifiée indépendamment
\end{itemize}

\textbf{Avantages :}
\begin{itemize}[leftmargin=*]
    \item Si une page est séparée du reste, elle reste vérifiable
    \item Pas besoin de chercher le QR code sur une page spécifique
    \item Meilleure traçabilité
\end{itemize}

\begin{tcolorbox}[colback=blue!5!white,colframe=blue!75!black,title=Logs de traitement multi-pages]
Dans la console de l'application, vous verrez :

\begin{verbatim}
📄 [QR-PDF] 2023001: 3 page(s) détectée(s)
  ├─ Ajout QR sur page 1/3
  ├─ Ajout QR sur page 2/3
  └─ Ajout QR sur page 3/3
✅ Document traité avec succès
\end{verbatim}
\end{tcolorbox}

\subsection{Différences entre orientations (relevé/attestation vs diplôme)}

\begin{table}[h]
\centering
\begin{tabular}{|l|c|c|}
\hline
\textbf{Critère} & \textbf{Relevé / Attestation} & \textbf{Diplôme} \\
\hline
Orientation & Portrait & Paysage \\
\hline
Dimensions A4 & 21 x 29.7 cm & 29.7 x 21 cm \\
\hline
Position "Haut Droite" & 85\%, 15\% & 90\%, 10\% \\
\hline
\end{tabular}
\caption{Différences selon le type de document}
\end{table}

\subsection{Tester la lisibilité d'un QR code}

Pour vérifier qu'un QR code généré est bien lisible :

\begin{enumerate}[leftmargin=*]
    \item Ouvrez le PDF sur votre ordinateur
    \item Utilisez un smartphone avec une application de scan de QR code
    \item Approchez la caméra du QR code à l'écran
    \item Le contenu devrait s'afficher immédiatement
    \item Vérifiez que toutes les informations sont correctes
\end{enumerate}

\textbf{Applications de scan recommandées :}
\begin{itemize}[leftmargin=*]
    \item iOS : Application Appareil photo native, QR Code Reader
    \item Android : Google Lens, QR Code Reader
\end{itemize}

\subsection{Limitation en mode démo}

\begin{tcolorbox}[colback=red!5!white,colframe=red!75!black,title=Restriction mode démo]
\textbf{Menu désactivé en mode démo}

Si vous n'avez pas de licence active :
\begin{itemize}[leftmargin=*]
    \item Le menu \menu{QR Codes sur PDF} apparaît grisé dans la barre latérale
    \item Un clic affiche : "Cette fonctionnalité nécessite une licence active"
    \item Vous devez activer l'application via \menu{Paramètres} > \menu{Licence}
\end{itemize}

\textbf{Alternative :} Utilisez le menu \menu{Placement QR sur Documents} qui reste disponible en mode démo.
\end{tcolorbox}

\subsection{Conseils pour un résultat optimal}

\begin{enumerate}[leftmargin=*]
    \item \textbf{Nommage des fichiers PDF} :
    \begin{itemize}
        \item Incluez le matricule dans le nom pour faciliter le matching automatique
        \item Format recommandé : \texttt{Type\_MATRICULE\_NOM\_Prenom.pdf}
        \item Exemple : \texttt{Releve\_2023001\_DUPONT\_Jean.pdf}
    \end{itemize}

    \item \textbf{Position du QR code} :
    \begin{itemize}
        \item Évitez le centre qui peut masquer des informations
        \item Préférez les coins (Haut Droite recommandé)
        \item Gardez une marge de 10\% par rapport aux bords
    \end{itemize}

    \item \textbf{Taille du QR code} :
    \begin{itemize}
        \item Trop petit : difficile à scanner
        \item Trop grand : envahissant visuellement
        \item Recommandation : Moyen (80x80 px) ou Grand (100x100 px)
    \end{itemize}

    \item \textbf{Contenu du QR code} :
    \begin{itemize}
        \item Limitez aux informations essentielles
        \item Maximum recommandé : 5-6 champs
        \item Plus le contenu est long, plus le QR est dense
    \end{itemize}

    \item \textbf{Correction d'erreur} :
    \begin{itemize}
        \item Utilisez "M" (Medium) par défaut
        \item Passez à "Q" ou "H" si les documents risquent d'être photocopiés plusieurs fois
    \end{itemize}
\end{enumerate}

\newpage

% Section 06 - Placement QR sur Documents
\section{Placement avancé de QR code sur documents}

\subsection{Vue d'ensemble}

Le menu \menu{Placement QR sur Documents} offre une fonctionnalité similaire à \menu{QR Codes sur PDF}, mais avec deux modes de fonctionnement : manuel et Excel.

\info{Cette fonctionnalité est DISPONIBLE même en mode démo, contrairement à "QR Codes sur PDF".}

\textbf{Principales différences avec "QR Codes sur PDF" :}

\begin{table}[h]
\centering
\begin{tabular}{|l|p{5cm}|p{5cm}|}
\hline
\textbf{Critère} & \textbf{QR Codes sur PDF} & \textbf{Placement QR sur Documents} \\
\hline
Mode démo & ✗ Désactivé & ✓ Disponible \\
\hline
Modes & Uniquement Excel & Manuel ET Excel \\
\hline
QR personnalisé & Basé sur colonnes Excel & Texte libre + Excel \\
\hline
Complexité & Simple & Avancée \\
\hline
\end{tabular}
\caption{Comparaison des deux menus QR}
\end{table}

\subsection{Mode 1 : Placement manuel (texte libre)}

\subsubsection{Cas d'usage}

Utilisez ce mode pour :
\begin{itemize}[leftmargin=*]
    \item Ajouter un QR code unique sur un seul document
    \item Créer un QR code avec du texte personnalisé (URL, vCard, texte libre)
    \item Tester rapidement le positionnement avant un traitement en masse
\end{itemize}

\subsubsection{Workflow manuel}

\begin{enumerate}[leftmargin=*]
    \item Cliquez sur \menu{Placement QR sur Documents}
    \item Sélectionnez le mode \textbf{Manuel}
    \item Chargez UN document PDF avec \bouton{Charger PDF}
    \item Dans le champ \champ{Contenu du QR code}, saisissez le texte souhaité :
    \begin{itemize}
        \item URL : \texttt{https://www.fmsp.edu/verify?id=123}
        \item Texte : \texttt{Document officiel FMSP 2024}
        \item vCard : Contact au format vCard
        \item Tout autre texte libre
    \end{itemize}
    \item Configurez la position avec la grille A4
    \item Choisissez la taille et la correction d'erreur
    \item Cliquez sur \bouton{Générer le QR code}
    \item Un aperçu s'affiche avec le QR code positionné
    \item Si satisfait, cliquez sur \bouton{Télécharger le PDF}
\end{enumerate}

\begin{figure}[h]
    \centering
    % \includegraphics[width=0.9\textwidth]{images/qr_placement_manuel.png}
    \fbox{\parbox{0.85\textwidth}{\centering [CAPTURE D'ÉCRAN : MODE MANUEL]\\ \vspace{0.3cm}
    Champ texte "Contenu du QR code"\\
    Grille de positionnement\\
    Aperçu du PDF avec QR code}}
    \caption{Interface du mode manuel}
\end{figure}

\subsection{Mode 2 : Traitement en lot avec Excel}

\subsubsection{Cas d'usage}

Utilisez ce mode pour :
\begin{itemize}[leftmargin=*]
    \item Traiter plusieurs documents en une seule opération
    \item Associer automatiquement des données Excel à chaque document
    \item Générer des QR codes uniques basés sur les données de chaque étudiant
\end{itemize}

\subsubsection{Workflow Excel}

\begin{enumerate}[leftmargin=*]
    \item Sélectionnez le mode \textbf{Excel}
    \item Chargez plusieurs PDFs avec \bouton{Charger Documents PDF}
    \item Chargez le fichier Excel avec \bouton{Charger Fichier Excel}
    \item Sélectionnez la \champ{Colonne de correspondance} :
    \begin{itemize}
        \item MATRICULE (recommandé)
        \item NOM
        \item Toute autre colonne unique
    \end{itemize}
    \item L'application tente de matcher automatiquement chaque PDF
    \item Corrigez manuellement les correspondances non trouvées
    \item Sélectionnez les colonnes à inclure dans les QR codes
    \item Configurez le positionnement et les paramètres du QR
    \item Choisissez le format d'export (ZIP / Individual / Single)
    \item Cliquez sur \bouton{Traiter tous les documents}
\end{enumerate}

\subsection{Paramètres avancés du QR code}

\subsubsection{Taille personnalisée}

\begin{enumerate}[leftmargin=*]
    \item Sélectionnez \champ{Taille : Personnalisée}
    \item Un champ s'affiche : \champ{Taille en pixels}
    \item Entrez une valeur entre 40 et 200 pixels
    \item Recommandations :
    \begin{itemize}
        \item 40-60 px : Très discret, pour documents chargés
        \item 80-100 px : Standard, bon compromis
        \item 120-150 px : Grand, très visible
        \item 150+ px : Très grand, pour affiches
    \end{itemize}
\end{enumerate}

\subsubsection{Correction d'erreur}

Les 4 niveaux expliqués en détail :

\begin{description}[leftmargin=2.5cm, style=nextline]
    \item[L (Low)]
    \begin{itemize}
        \item 7\% de données récupérables si endommagées
        \item QR code plus simple (moins de modules)
        \item Idéal pour : documents imprimés sur imprimante de qualité, environnement propre
        \item Éviter si : photocopies multiples, scan de mauvaise qualité
    \end{itemize}

    \item[M (Medium)]
    \begin{itemize}
        \item 15\% de données récupérables
        \item Bon équilibre entre densité et robustesse
        \item \textbf{Recommandé pour la plupart des cas}
        \item Usage type : relevés de notes, attestations standard
    \end{itemize}

    \item[Q (Quartile)]
    \begin{itemize}
        \item 25\% de données récupérables
        \item QR code plus dense
        \item Idéal pour : documents qui seront photocopiés, archivage long terme
        \item Usage type : diplômes, certificats officiels
    \end{itemize}

    \item[H (High)]
    \begin{itemize}
        \item 30\% de données récupérables
        \item QR code très dense
        \item Idéal pour : environnements difficiles, affichage public, documents manipulés fréquemment
        \item Usage type : affiches, documents d'identité
    \end{itemize}
\end{description}

\subsection{Prévisualisation avant traitement}

Avant de lancer un traitement en masse :

\begin{enumerate}[leftmargin=*]
    \item Configurez tous les paramètres
    \item Cliquez sur \bouton{Prévisualiser avec le premier document}
    \item Un PDF de test s'affiche avec le QR code positionné
    \item Vérifiez :
    \begin{itemize}
        \item Position correcte
        \item Taille appropriée
        \item Lisibilité du QR code (testez avec un scanner)
        \item Absence de chevauchement avec du texte important
    \end{itemize}
    \item Si nécessaire, ajustez les paramètres et prévisualisez à nouveau
    \item Une fois satisfait, lancez le traitement complet
\end{enumerate}

\subsection{Gestion des correspondances}

\subsubsection{Matching automatique}

L'algorithme de matching utilise plusieurs stratégies :

\begin{enumerate}[leftmargin=*]
    \item \textbf{Recherche exacte} : Le nom du fichier contient la valeur exacte de la colonne de correspondance
    \item \textbf{Recherche partielle} : Recherche insensible à la casse et aux accents
    \item \textbf{Extraction de mots-clés} : Si le fichier est nommé "Releve\_2023001\_DUPONT.pdf", extraction de "2023001" et "DUPONT"
\end{enumerate}

\subsubsection{Correction manuelle}

Pour les documents non matchés automatiquement :

\begin{enumerate}[leftmargin=*]
    \item Le document apparaît en \textcolor{red}{rouge} avec statut "Non matché"
    \item Cliquez sur \bouton{Sélectionner manuellement}
    \item Une liste déroulante s'affiche avec tous les étudiants de l'Excel
    \item Utilisez la barre de recherche pour trouver rapidement
    \item Sélectionnez l'étudiant correspondant
    \item Le statut passe à \textcolor{green}{✓ Matché}
\end{enumerate}

\subsection{Formats d'exportation}

Les mêmes 3 formats que les autres fonctionnalités :

\begin{description}[leftmargin=3cm]
    \item[ZIP]
    \begin{itemize}
        \item Tous les documents dans une archive
        \item Nom : \texttt{Documents\_QR\_[Date]\_[Heure].zip}
    \end{itemize}

    \item[Fichiers individuels]
    \begin{itemize}
        \item Chaque document téléchargé séparément
        \item Délai de 100 ms entre chaque téléchargement
        \item Nom conservé : \texttt{[Nom\_original]\_avec\_QR.pdf}
    \end{itemize}

    \item[PDF unique]
    \begin{itemize}
        \item Tous les documents fusionnés
        \item Nom : \texttt{Documents\_Fusionnes\_QR\_[Date].pdf}
        \item Ordre : alphabétique par nom de fichier
    \end{itemize}
\end{description}

\newpage

% Section 07 - Configuration académique
\section{Configuration académique des classes}

\subsection{Menu "Config des relevés"}

Ce menu permet de créer et gérer les configurations de classes nécessaires pour générer les relevés de notes.

\subsection{Structure hiérarchique}

Une configuration de classe est organisée en 3 niveaux :

\begin{verbatim}
Configuration de Classe
├── Semestre 1
│   ├── UE1 - Sciences Fondamentales (30 crédits)
│   │   ├── EC: Mathématiques (poids: 2, base: 20)
│   │   ├── EC: Physique (poids: 1.5, base: 20)
│   │   └── EC: Chimie (poids: 2, base: 20)
│   └── UE2 - Sciences Humaines (30 crédits)
│       └── EC: Anglais (poids: 1, base: 20)
├── Semestre 2
│   └── ...
└── Semestre 3
    └── ...
\end{verbatim}

\subsection{Créer une nouvelle configuration}

\subsubsection{Étape 1 : Informations générales}

\begin{enumerate}[leftmargin=*]
    \item Cliquez sur \bouton{Nouvelle Configuration}
    \item Remplissez :
    \begin{itemize}
        \item \champ{Nom de la configuration} : Ex. "Licence 1 Médecine - 2023-2024"
        \item \champ{Niveau} : L1, L2, L3, M1, M2, etc.
    \end{itemize}
    \item Cliquez sur \bouton{Créer}
\end{enumerate}

\subsubsection{Étape 2 : Ajouter des semestres}

\begin{enumerate}[leftmargin=*]
    \item Dans la configuration créée, cliquez sur \bouton{Ajouter un semestre}
    \item Remplissez :
    \begin{itemize}
        \item \champ{Nom du semestre} : "Semestre 1" ou "S1"
        \item \champ{Code} : S1, S2, S3, etc.
    \end{itemize}
    \item Répétez pour ajouter tous les semestres nécessaires
\end{enumerate}

\subsubsection{Étape 3 : Ajouter des UEs}

Pour chaque semestre :

\begin{enumerate}[leftmargin=*]
    \item Cliquez sur \bouton{Ajouter une UE}
    \item Remplissez :
    \begin{itemize}
        \item \champ{Nom de l'UE} : "UE1 - Sciences Fondamentales"
        \item \champ{Crédits ECTS} : Nombre de crédits (ex: 30)
    \end{itemize}
    \item Cliquez sur \bouton{Créer l'UE}
\end{enumerate}

\subsubsection{Étape 4 : Ajouter des ECs}

Pour chaque UE :

\begin{enumerate}[leftmargin=*]
    \item Cliquez sur \bouton{Ajouter un EC}
    \item Remplissez :
    \begin{itemize}
        \item \champ{Nom de l'EC} : "Mathématiques", "Physique", etc.
        \item \champ{Poids (coefficient)} : 1, 1.5, 2, etc.
        \item \champ{Base de notation} : 10, 20, ou 100
        \item \champ{A une session de rattrapage} : Cochez si applicable
    \end{itemize}
    \item Cliquez sur \bouton{Créer l'EC}
    \item Répétez pour tous les ECs de l'UE
\end{enumerate}

\begin{figure}[h]
    \centering
    % \includegraphics[width=\textwidth]{images/config_classe.png}
    \fbox{\parbox{0.9\textwidth}{\centering [CAPTURE D'ÉCRAN : INTERFACE DE CONFIGURATION]\\ \vspace{0.3cm}
    Arborescence : Config > Semestre > UE > EC\\
    Boutons "Ajouter" à chaque niveau\\
    Formulaires d'édition}}
    \caption{Interface de configuration académique}
\end{figure}

\subsection{Réorganiser les UEs}

Par défaut, les UEs s'affichent dans l'ordre de création. Pour les réorganiser :

\begin{enumerate}[leftmargin=*]
    \item Cliquez sur \bouton{Réorganiser les UEs}
    \item Utilisez le drag \& drop pour réordonner
    \item Ou utilisez les boutons ↑ et ↓
    \item Cliquez sur \bouton{Enregistrer l'ordre}
\end{enumerate}

\subsection{Personnaliser le thème des relevés}

Chaque configuration peut avoir son propre thème :

\begin{enumerate}[leftmargin=*]
    \item Dans la configuration, cliquez sur l'onglet \textbf{Thème}
    \item Configurez :
    \begin{itemize}
        \item Couleurs (primaire, secondaire, accent)
        \item Polices (famille, tailles)
        \item Mise en page
    \end{itemize}
    \item Prévisualisez en temps réel
    \item Enregistrez le thème
\end{enumerate}

\subsection{Exporter / Importer une configuration}

\subsubsection{Exporter}

\begin{enumerate}[leftmargin=*]
    \item Sélectionnez la configuration
    \item Cliquez sur \bouton{Exporter la configuration}
    \item Un fichier JSON est téléchargé : \texttt{config\_[Nom].json}
\end{enumerate}

\subsubsection{Importer}

\begin{enumerate}[leftmargin=*]
    \item Cliquez sur \bouton{Importer une configuration}
    \item Sélectionnez le fichier JSON
    \item La configuration est chargée et ajoutée à votre liste
\end{enumerate}

\astuce{Partagez vos configurations avec vos collègues via export/import pour standardiser les structures académiques.}

\subsection{Menu "Configurer les entêtes"}

Ce menu permet de définir les informations de l'établissement qui apparaîtront sur tous les documents.

\subsubsection{Informations de base}

\begin{itemize}[leftmargin=*]
    \item \champ{Nom de l'établissement}
    \item \champ{Adresse complète}
    \item \champ{Téléphone}
    \item \champ{Email}
    \item \champ{Site web}
\end{itemize}

\subsubsection{Logo}

\begin{enumerate}[leftmargin=*]
    \item Cliquez sur \bouton{Charger un logo}
    \item Sélectionnez une image PNG ou JPG
    \item Recommandations :
    \begin{itemize}
        \item Résolution : 800x200 pixels minimum
        \item Taille : Maximum 2 Mo
        \item Format : PNG avec transparence recommandé
    \end{itemize}
    \item Le logo s'affichera en haut de tous les relevés et attestations
\end{enumerate}

\subsection{Menu "Historique des docs"}

Ce menu permet de consulter l'historique de tous les documents générés.

\subsubsection{Informations affichées}

\begin{itemize}[leftmargin=*]
    \item Date et heure de génération
    \item Type de document (Relevé / Attestation)
    \item Nombre de documents générés
    \item Configuration utilisée
    \item Semestre(s) concerné(s)
    \item Format d'export utilisé
\end{itemize}

\subsubsection{Recherche et filtrage}

\begin{enumerate}[leftmargin=*]
    \item Utilisez la barre de recherche pour trouver une génération spécifique
    \item Filtrez par :
    \begin{itemize}
        \item Type de document
        \item Date (période)
        \item Configuration
    \end{itemize}
    \item Cliquez sur une entrée pour voir les détails complets
\end{enumerate}

\subsubsection{Export de l'historique}

\begin{enumerate}[leftmargin=*]
    \item Cliquez sur \bouton{Exporter l'historique}
    \item Choisissez le format : Excel ou PDF
    \item Un fichier récapitulatif est téléchargé
\end{enumerate}

\newpage

% Section 08 - Dépannage et résolution de problèmes
\section{Dépannage et résolution de problèmes}

\subsection{Problèmes courants et solutions}

\subsubsection{L'application ne se lance pas}

\textbf{Symptômes :}
\begin{itemize}[leftmargin=*]
    \item Double-clic sur l'icône sans effet
    \item Fenêtre s'affiche brièvement puis disparaît
    \item Message d'erreur au démarrage
\end{itemize}

\textbf{Solutions :}
\begin{enumerate}[leftmargin=*]
    \item \textbf{Vérifier la configuration système}
    \begin{itemize}
        \item RAM disponible : minimum 4 Go
        \item Espace disque : minimum 500 Mo
        \item Système d'exploitation à jour
    \end{itemize}

    \item \textbf{Désinstaller et réinstaller}
    \begin{itemize}
        \item Désinstallez complètement l'application
        \item Téléchargez la dernière version
        \item Réinstallez
    \end{itemize}

    \item \textbf{Vérifier les permissions}
    \begin{itemize}
        \item Windows : Exécuter en tant qu'administrateur
        \item macOS : Autoriser dans Sécurité et confidentialité
        \item Linux : Vérifier les permissions d'exécution
    \end{itemize}
\end{enumerate}

\subsubsection{Erreur lors du chargement du fichier Excel}

\textbf{Message d'erreur :} "Impossible de lire le fichier Excel" ou "Colonnes manquantes"

\textbf{Solutions :}
\begin{enumerate}[leftmargin=*]
    \item \textbf{Vérifier le format du fichier}
    \begin{itemize}
        \item Format accepté : .xlsx ou .xls
        \item Pas de macros actives
        \item Fichier non protégé par mot de passe
    \end{itemize}

    \item \textbf{Vérifier les colonnes obligatoires}
    \begin{itemize}
        \item Pour relevés : MATRICULE, NOM, PRENOM, SEMESTRE, NIVEAU, FILIERE
        \item Pour attestations : Ajouter DATE\_NAISSANCE, LIEU\_NAISSANCE, DOMAINE, PARCOURS, MOYENNE\_GENERALE
        \item Respect de la casse (majuscules/minuscules)
    \end{itemize}

    \item \textbf{Fermer Excel}
    \begin{itemize}
        \item Le fichier ne doit pas être ouvert dans Excel
        \item Fermez Excel complètement
        \item Réessayez le chargement
    \end{itemize}

    \item \textbf{Vérifier l'encodage}
    \begin{itemize}
        \item Enregistrez le fichier avec encodage UTF-8
        \item Évitez les caractères spéciaux dans les noms de colonnes
    \end{itemize}
\end{enumerate}

\subsubsection{Erreur "Out of Memory" (OOM)}

\textbf{Message :} "JavaScript heap out of memory" ou "Mémoire insuffisante"

\textbf{Causes :}
\begin{itemize}[leftmargin=*]
    \item Génération de trop de documents simultanément
    \item Fichier Excel très volumineux
    \item Autres applications consommant beaucoup de RAM
\end{itemize}

\textbf{Solutions :}
\begin{enumerate}[leftmargin=*]
    \item \textbf{Réduire le nombre de documents}
    \begin{itemize}
        \item Limitez à 200 relevés par génération
        \item Pour plus, divisez en plusieurs lots
        \item Utilisez les filtres pour sélectionner un sous-ensemble
    \end{itemize}

    \item \textbf{Fermer les autres applications}
    \begin{itemize}
        \item Libérez de la RAM
        \item Fermez les navigateurs avec beaucoup d'onglets
        \item Redémarrez l'ordinateur si nécessaire
    \end{itemize}

    \item \textbf{Simplifier le fichier Excel}
    \begin{itemize}
        \item Supprimez les colonnes inutiles
        \item Retirez les formules complexes
        \item Conservez uniquement les données nécessaires
    \end{itemize}
\end{enumerate}

\subsubsection{Mapping des colonnes non sauvegardé}

\textbf{Symptôme :} À chaque chargement d'Excel, le mapping doit être refait

\textbf{Cause :} Changement des noms de colonnes entre les chargements

\textbf{Solution :}
\begin{itemize}[leftmargin=*]
    \item Gardez les mêmes noms de colonnes dans vos fichiers Excel
    \item Le mapping est basé sur le NOM de la colonne, pas sa position
    \item Si vous devez renommer, utilisez \bouton{Réinitialiser le mapping}
\end{itemize}

\subsubsection{QR codes non lisibles}

\textbf{Symptômes :}
\begin{itemize}[leftmargin=*]
    \item Le scanner ne détecte pas le QR code
    \item Message d'erreur lors du scan
    \item QR code trop petit ou trop flou
\end{itemize}

\textbf{Solutions :}
\begin{enumerate}[leftmargin=*]
    \item \textbf{Augmenter la taille}
    \begin{itemize}
        \item Passez de "Petit" à "Moyen" ou "Grand"
        \item Minimum recommandé : 80x80 pixels
    \end{itemize}

    \item \textbf{Augmenter la correction d'erreur}
    \begin{itemize}
        \item Passez de "L" ou "M" à "Q" ou "H"
        \item Plus de redondance = meilleure lisibilité
    \end{itemize}

    \item \textbf{Réduire le contenu}
    \begin{itemize}
        \item Limitez le nombre de colonnes dans le QR
        \item Trop de données = QR trop dense
        \item Maximum recommandé : 5-6 champs
    \end{itemize}

    \item \textbf{Vérifier la position}
    \begin{itemize}
        \item Le QR ne doit pas chevaucher du texte
        \item Fond blanc recommandé autour du QR
        \item Évitez les zones avec images de fond
    \end{itemize}
\end{enumerate}

\subsubsection{Moyennes calculées incorrectement}

\textbf{Symptôme :} Les moyennes affichées ne correspondent pas aux calculs manuels

\textbf{Vérifications :}
\begin{enumerate}[leftmargin=*]
    \item \textbf{Poids des ECs}
    \begin{itemize}
        \item Vérifiez les coefficients dans la configuration
        \item Formule : $\text{Moy UE} = \frac{\sum(\text{Note} \times \text{Poids})}{\sum \text{Poids}}$
    \end{itemize}

    \item \textbf{Crédits des UEs}
    \begin{itemize}
        \item Les crédits ECTS sont-ils corrects ?
        \item Formule : $\text{Moy Gén} = \frac{\sum(\text{Moy UE} \times \text{Crédits})}{\sum \text{Crédits}}$
    \end{itemize}

    \item \textbf{Notes manquantes}
    \begin{itemize}
        \item Les cellules vides sont ignorées dans le calcul
        \item Assurez-vous que toutes les notes sont présentes
    \end{itemize}

    \item \textbf{Base de notation}
    \begin{itemize}
        \item Vérifiez que la base définie (10, 20, 100) correspond aux notes
        \item L'application normalise automatiquement sur 20
    \end{itemize}
\end{enumerate}

\subsubsection{Attestations générées pour des étudiants ajournés}

\textbf{Symptôme :} Des étudiants avec moyenne < 10 reçoivent des attestations

\attention{Ceci ne devrait PAS se produire. L'application filtre automatiquement.}

\textbf{Vérifications :}
\begin{enumerate}[leftmargin=*]
    \item Vérifiez la colonne \texttt{MOYENNE\_GENERALE} dans l'Excel
    \item Les valeurs doivent être numériques (pas de texte)
    \item Format : 10.5, 12.75, etc. (point décimal, pas virgule)
    \item Si le problème persiste, contactez le support
\end{enumerate}

\subsubsection{Thème personnalisé non appliqué}

\textbf{Symptôme :} Le thème configuré n'apparaît pas sur les documents générés

\textbf{Solutions :}
\begin{enumerate}[leftmargin=*]
    \item Vérifiez que le thème est bien enregistré (bouton \bouton{Enregistrer le thème})
    \item Pour les relevés : le thème est lié à la configuration de classe
    \item Pour les attestations : le thème est indépendant
    \item Essayez de prévisualiser pour voir si le thème s'applique
    \item Exportez et réimportez le thème pour tester
\end{enumerate}

\subsection{Optimisation des performances}

\subsubsection{Accélérer la génération}

\begin{enumerate}[leftmargin=*]
    \item \textbf{Désactiver la compression ZIP} (si non nécessaire)
    \begin{itemize}
        \item Gain de temps : 20-30\%
        \item Fichiers plus volumineux
    \end{itemize}

    \item \textbf{Générer par lots de 100-200}
    \begin{itemize}
        \item Au lieu de 500 en une fois
        \item Évite les erreurs mémoire
        \item Permet de vérifier progressivement
    \end{itemize}

    \item \textbf{Simplifier les thèmes}
    \begin{itemize}
        \item Évitez trop d'images ou de polices complexes
        \item Utilisez des préréglages plutôt que des thèmes très personnalisés
    \end{itemize}

    \item \textbf{Utiliser un SSD}
    \begin{itemize}
        \item Plus rapide qu'un disque dur mécanique
        \item Améliore l'écriture des fichiers générés
    \end{itemize}
\end{enumerate}

\subsubsection{Réduire l'utilisation mémoire}

\begin{enumerate}[leftmargin=*]
    \item Fermez toutes les applications inutiles
    \item Utilisez le format "Fichiers individuels" au lieu de "ZIP" si possible
    \item Ne générez pas de QR codes si non nécessaire
    \item Redémarrez l'application entre deux générations massives
\end{enumerate}

\subsection{Journalisation et logs}

\subsubsection{Consulter les logs de l'application}

Les logs peuvent aider à identifier des problèmes :

\textbf{Windows :}
\begin{verbatim}
%APPDATA%\FMSP-Releves\logs\
\end{verbatim}

\textbf{macOS :}
\begin{verbatim}
~/Library/Logs/FMSP-Releves/
\end{verbatim}

\textbf{Linux :}
\begin{verbatim}
~/.config/FMSP-Releves/logs/
\end{verbatim}

\subsubsection{Console de développement}

Pour les utilisateurs avancés :

\begin{enumerate}[leftmargin=*]
    \item Appuyez sur \texttt{F12} ou \texttt{Ctrl+Shift+I} (Windows/Linux)
    \item Ou \texttt{Cmd+Option+I} (macOS)
    \item La console de développement s'ouvre
    \item Consultez l'onglet \textbf{Console} pour les messages
    \item Recherchez les messages en rouge (erreurs)
\end{enumerate}

\subsection{Réinitialisation de l'application}

Si les problèmes persistent, réinitialisez l'application :

\begin{tcolorbox}[colback=red!5!white,colframe=red!75!black,title=ATTENTION : Sauvegardez vos données]
La réinitialisation supprime :
\begin{itemize}[leftmargin=*]
    \item Toutes les configurations de classes
    \item Tous les thèmes personnalisés
    \item L'historique des générations
    \item Les paramètres de l'établissement
\end{itemize}

Exportez vos configurations et thèmes avant de continuer !
\end{tcolorbox}

\textbf{Procédure :}
\begin{enumerate}[leftmargin=*]
    \item Menu \menu{Paramètres} > \textbf{Avancé}
    \item Cliquez sur \bouton{Réinitialiser l'application}
    \item Confirmez l'action
    \item L'application redémarre avec les paramètres par défaut
\end{enumerate}

\subsection{Contacter le support}

Si aucune solution ne fonctionne :

\begin{enumerate}[leftmargin=*]
    \item Menu \menu{Aide} > \bouton{Rapporter un problème}
    \item Remplissez le formulaire :
    \begin{itemize}
        \item Description du problème
        \item Étapes pour reproduire
        \item Captures d'écran si possible
    \end{itemize}
    \item Joignez les logs récents
    \item Cliquez sur \bouton{Envoyer}
\end{enumerate}

\textbf{Informations utiles à fournir :}
\begin{itemize}[leftmargin=*]
    \item Version de l'application (Menu \menu{Aide} > \menu{À propos})
    \item Système d'exploitation et version
    \item Nombre d'étudiants traités
    \item Type de document généré
    \item Message d'erreur exact (copier-coller)
\end{itemize}

\newpage

% Section 09 - Cas d'usage pratiques
\section{Cas d'usage pratiques et scénarios}

\subsection{Scénario 1 : Fin de semestre - Relevés pour une classe de 150 étudiants}

\textbf{Contexte :}
\begin{itemize}[leftmargin=*]
    \item Classe : Licence 1 Médecine
    \item Semestre : S1
    \item Nombre d'étudiants : 150
    \item Objectif : Générer tous les relevés pour distribution
\end{itemize}

\textbf{Workflow détaillé :}

\begin{enumerate}[leftmargin=*]
    \item \textbf{Préparation} (15 min)
    \begin{itemize}
        \item Créer la configuration "L1 Médecine 2023-2024" dans \menu{Config des relevés}
        \item Ajouter le semestre S1 avec toutes les UEs et ECs
        \item Configurer les coefficients et crédits
        \item Personnaliser le thème (optionnel)
    \end{itemize}

    \item \textbf{Préparation du fichier Excel} (10 min)
    \begin{itemize}
        \item Vérifier toutes les colonnes obligatoires présentes
        \item Vérifier que les noms d'ECs correspondent exactement à la configuration
        \item Enregistrer au format .xlsx
    \end{itemize}

    \item \textbf{Génération} (5-10 min)
    \begin{itemize}
        \item Charger le fichier Excel (150 lignes)
        \item Sélectionner la configuration "L1 Médecine 2023-2024"
        \item Choisir le semestre S1
        \item Mapper les colonnes (mémorisé pour la prochaine fois)
        \item Format d'export : ZIP avec compression
        \item QR code chiffré : Activé
        \item Lancer la génération
    \end{itemize}

    \item \textbf{Vérification} (5 min)
    \begin{itemize}
        \item Ouvrir 5-10 relevés au hasard
        \item Vérifier les notes, moyennes, mentions
        \item Tester un QR code
    \end{itemize}

    \item \textbf{Distribution}
    \begin{itemize}
        \item Extraire le ZIP (150 fichiers PDF)
        \item Imprimer ou distribuer numériquement
    \end{itemize}
\end{enumerate}

\textbf{Résultat attendu :}
\begin{itemize}[leftmargin=*]
    \item 150 relevés PDF individuels dans une archive ZIP
    \item Taille totale : environ 30-40 Mo (avec compression)
    \item Temps total : 35-40 minutes
    \item Chaque relevé contient toutes les notes, moyennes calculées, mention, et QR code
\end{itemize}

\subsection{Scénario 2 : Génération multi-semestres pour relevé de parcours}

\textbf{Contexte :}
\begin{itemize}[leftmargin=*]
    \item Étudiant finissant sa Licence (L3)
    \item Besoin : Relevé complet de L1 à L3 (6 semestres)
    \item Pour : Dossier de candidature en Master
\end{itemize}

\textbf{Workflow :}

\begin{enumerate}[leftmargin=*]
    \item Charger le fichier Excel contenant TOUS les semestres
    \item Sélectionner la configuration de classe
    \item Choisir \textbf{"Tous les semestres"}
    \item Cocher S1, S2, S3, S4, S5, S6
    \item Format d'export : \textbf{PDF unique} (tout fusionné en un fichier)
    \item Générer
\end{enumerate}

\textbf{Résultat :}
\begin{itemize}[leftmargin=*]
    \item 1 fichier PDF de 6 pages (1 page par semestre)
    \item Nom : \texttt{Releves\_Fusionnes\_DUPONT\_Jean\_2024.pdf}
    \item Parfait pour candidatures et archivage
\end{itemize}

\subsection{Scénario 3 : Attestations de réussite en fin d'année}

\textbf{Contexte :}
\begin{itemize}[leftmargin=*]
    \item Fin de l'année académique
    \item 200 étudiants de L1
    \item Besoin : Attestations pour les étudiants admis uniquement
\end{itemize}

\textbf{Workflow :}

\begin{enumerate}[leftmargin=*]
    \item Menu \menu{Générer les attestations}
    \item Charger le fichier Excel avec toutes les notes et \texttt{MOYENNE\_GENERALE}
    \item L'application affiche automatiquement :
    \begin{itemize}
        \item 167 étudiants ADMIS (moyenne $\geq$ 10)
        \item 33 étudiants AJOURNÉS (exclus automatiquement)
        \item Taux de réussite : 83.5\%
    \end{itemize}
    \item Choisir un préréglage de thème : "Professionnel Académique"
    \item Activer le QR code chiffré
    \item Format : ZIP avec compression
    \item Générer
\end{enumerate}

\textbf{Résultat :}
\begin{itemize}[leftmargin=*]
    \item 167 attestations PDF (uniquement les admis)
    \item Archive ZIP de 15-20 Mo
    \item Temps de génération : 8-10 minutes
\end{itemize}

\subsection{Scénario 4 : Ajout de QR codes sur relevés déjà imprimés}

\textbf{Contexte :}
\begin{itemize}[leftmargin=*]
    \item 100 relevés déjà générés et distribués SANS QR code
    \item Nouvelle politique : Tous les relevés doivent avoir un QR de vérification
    \item Documents déjà scannés en PDF
\end{itemize}

\textbf{Workflow :}

\begin{enumerate}[leftmargin=*]
    \item Menu \menu{QR Codes sur PDF} (nécessite licence active)
    \item Charger les 100 PDFs scannés
    \item Charger le fichier Excel avec les données étudiants
    \item Vérifier le matching automatique (devrait être 100\% si noms de fichiers contiennent matricules)
    \item Corriger manuellement les non-matchés
    \item Positionner le QR : Haut Droite (85\%, 15\%)
    \item Colonnes dans QR : MATRICULE, NOM, PRENOM, SEMESTRE
    \item Taille : Moyen, Correction d'erreur : M
    \item Format : ZIP
    \item Traiter tous les documents
\end{enumerate}

\textbf{Résultat :}
\begin{itemize}[leftmargin=*]
    \item 100 PDFs avec QR codes ajoutés
    \item QR présent sur toutes les pages de chaque document
    \item Prêts à être réimprimés
\end{itemize}

\subsection{Scénario 5 : Création d'attestations bilingues (FR/EN)}

\textbf{Contexte :}
\begin{itemize}[leftmargin=*]
    \item Étudiants internationaux demandant des attestations en anglais
    \item 50 étudiants admis
    \item Besoin : Attestations en anglais
\end{itemize}

\textbf{Préparation du fichier Excel :}

Ajouter les colonnes de traduction :
\begin{itemize}[leftmargin=*]
    \item \texttt{DOMAINE\_EN} : "Health Sciences"
    \item \texttt{PARCOURS\_EN} : "General Medicine"
    \item \texttt{SPECIALITE\_EN} : "Cardiology"
    \item \texttt{MENTION\_EN} : "Excellent", "Good", "Fair", "Pass"
\end{itemize}

\textbf{Workflow :}

\begin{enumerate}[leftmargin=*]
    \item Charger l'Excel avec colonnes FR et EN
    \item Vérifier l'éligibilité (50 admis détectés)
    \item Dans l'onglet \textbf{Paramètres}, sélectionner \champ{Langue : Anglais}
    \item Choisir un thème approprié
    \item Générer
\end{enumerate}

\textbf{Résultat :}
\begin{itemize}[leftmargin=*]
    \item 50 attestations en anglais
    \item Titre : "Certificate of Achievement"
    \item Toutes les valeurs traduites
\end{itemize}

\subsection{Scénario 6 : Test et configuration initiale}

\textbf{Contexte :}
\begin{itemize}[leftmargin=*]
    \item Première utilisation de l'application
    \item Besoin : Configurer et tester avant utilisation en production
\end{itemize}

\textbf{Workflow recommandé :}

\begin{enumerate}[leftmargin=*]
    \item \textbf{Jour 1 : Configuration de base}
    \begin{itemize}
        \item Menu \menu{Configurer les entêtes}
        \item Renseigner toutes les informations de l'établissement
        \item Charger le logo
        \item Enregistrer
    \end{itemize}

    \item \textbf{Jour 2 : Créer une configuration de test}
    \begin{itemize}
        \item Menu \menu{Config des relevés}
        \item Créer une configuration simplifiée : "Test - L1"
        \item 1 semestre, 2 UEs, 3-4 ECs par UE
        \item Enregistrer
    \end{itemize}

    \item \textbf{Jour 3 : Test avec données fictives}
    \begin{itemize}
        \item Créer un fichier Excel avec 5-10 lignes de données fictives
        \item Générer des relevés de test
        \item Vérifier le résultat
        \item Ajuster le thème si nécessaire
    \end{itemize}

    \item \textbf{Jour 4 : Test d'attestations}
    \begin{itemize}
        \item Ajouter colonne \texttt{MOYENNE\_GENERALE} à l'Excel de test
        \item Générer des attestations de test
        \item Tester les différents thèmes prédéfinis
        \item Choisir celui qui convient
    \end{itemize}

    \item \textbf{Jour 5 : Configuration réelle}
    \begin{itemize}
        \item Créer les vraies configurations (toutes les classes)
        \item Personnaliser les thèmes définitifs
        \item Exporter les configurations pour sauvegarde
        \item Prêt pour production
    \end{itemize}
\end{enumerate}

\subsection{Scénario 7 : Gestion d'une promotion complète (L1 à L3)}

\textbf{Contexte :}
\begin{itemize}[leftmargin=*]
    \item Faculté avec 3 niveaux (L1, L2, L3)
    \item Environ 500 étudiants au total
    \item Fin de semestre 2
\end{itemize}

\textbf{Stratégie organisée :}

\begin{enumerate}[leftmargin=*]
    \item \textbf{Créer 3 configurations séparées}
    \begin{itemize}
        \item Config L1 (S1 et S2)
        \item Config L2 (S3 et S4)
        \item Config L3 (S5 et S6)
    \end{itemize}

    \item \textbf{Préparer 3 fichiers Excel}
    \begin{itemize}
        \item Excel\_L1.xlsx (150 étudiants)
        \item Excel\_L2.xlsx (180 étudiants)
        \item Excel\_L3.xlsx (170 étudiants)
    \end{itemize}

    \item \textbf{Générer par niveau}
    \begin{itemize}
        \item Jour 1 : Relevés L1 pour S2
        \item Jour 2 : Relevés L2 pour S4
        \item Jour 3 : Relevés L3 pour S6
        \item Vérification quotidienne des résultats
    \end{itemize}

    \item \textbf{Attestations pour les finissants L3}
    \begin{itemize}
        \item Générer attestations uniquement pour L3 admis
        \item Thème plus solennel
        \item QR code obligatoire
    \end{itemize}
\end{enumerate}

\textbf{Avantages de cette approche :}
\begin{itemize}[leftmargin=*]
    \item Évite la surcharge mémoire
    \item Permet des vérifications intermédiaires
    \item Facilite la correction d'erreurs
    \item Organisation claire par niveau
\end{itemize}

\subsection{Conseils généraux pour l'utilisation quotidienne}

\begin{enumerate}[leftmargin=*]
    \item \textbf{Gardez vos fichiers Excel organisés}
    \begin{itemize}
        \item Nommage cohérent : \texttt{Notes\_L1\_S1\_2023-2024.xlsx}
        \item Dossier dédié par année académique
        \item Sauvegardes régulières
    \end{itemize}

    \item \textbf{Exportez vos configurations régulièrement}
    \begin{itemize}
        \item Sauvegarde mensuelle des configurations
        \item Stockage sur cloud ou disque externe
        \item Facilite le partage entre collègues
    \end{itemize}

    \item \textbf{Testez toujours avec un petit échantillon}
    \begin{itemize}
        \item Avant une génération massive, testez avec 5-10 étudiants
        \item Vérifiez le résultat
        \item Puis lancez la génération complète
    \end{itemize}

    \item \textbf{Utilisez l'historique}
    \begin{itemize}
        \item Consultez régulièrement \menu{Historique des docs}
        \item Tracez toutes vos générations
        \item Exportez l'historique trimestriellement
    \end{itemize}

    \item \textbf{Standardisez vos thèmes}
    \begin{itemize}
        \item Créez un thème officiel pour chaque type de document
        \item Exportez-les
        \item Partagez avec vos collègues pour cohérence
    \end{itemize}
\end{enumerate}

\newpage

% Section 10 - Annexes
\section{Annexes}

\subsection{Raccourcis clavier}

L'application propose plusieurs raccourcis clavier pour accélérer le travail.

\subsubsection{Raccourcis généraux}

\begin{longtable}{|p{4cm}|p{10cm}|}
\hline
\textbf{Raccourci} & \textbf{Action} \\
\hline
\endfirsthead
\hline
\textbf{Raccourci} & \textbf{Action} \\
\hline
\endhead
\hline
\endfoot

Ctrl + N & Créer une nouvelle configuration de classe \\
\hline
Ctrl + O & Ouvrir une configuration existante \\
\hline
Ctrl + S & Sauvegarder la configuration en cours \\
\hline
Ctrl + E & Charger un fichier Excel \\
\hline
Ctrl + Q & Quitter l'application \\
\hline
F5 & Actualiser l'interface \\
\hline
F1 & Afficher l'aide \\
\hline
F12 & Ouvrir la console de développement (avancé) \\
\hline
\end{longtable}

\subsubsection{Raccourcis spécifiques}

\begin{longtable}{|p{4cm}|p{10cm}|}
\hline
\textbf{Raccourci} & \textbf{Action} \\
\hline
\endfirsthead
\hline
\textbf{Raccourci} & \textbf{Action} \\
\hline
\endhead
\hline
\endfoot

Ctrl + G & Générer les relevés de notes \\
\hline
Ctrl + A & Générer les attestations \\
\hline
Ctrl + P & Prévisualiser un document \\
\hline
Ctrl + Shift + E & Export de la configuration \\
\hline
Ctrl + Shift + I & Import de configuration \\
\hline
\end{longtable}

\subsection{Format des codes QR}

\subsubsection{Contenu standard}

Les codes QR générés contiennent les informations au format texte brut :

\begin{verbatim}
MATRICULE: 2023001
NOM: DUPONT
PRENOM: Jean
SEMESTRE: S1
NIVEAU: L1
FILIERE: MEDECINE
\end{verbatim}

\subsubsection{Chiffrement}

Lorsque le chiffrement est activé :

\begin{itemize}[leftmargin=*]
    \item \textbf{Algorithme} : AES-128-ECB
    \item \textbf{Clé} : Basée sur le matricule de l'étudiant
    \item \textbf{Format} : Texte chiffré en Base64
    \item \textbf{Avantage} : Impossible de falsifier sans la clé
\end{itemize}

\subsubsection{Niveaux de correction d'erreur}

\begin{table}[h]
\centering
\begin{tabular}{|c|c|p{8cm}|}
\hline
\textbf{Niveau} & \textbf{Récupération} & \textbf{Usage recommandé} \\
\hline
L & 7\% & Documents de qualité, impression professionnelle \\
\hline
M & 15\% & Usage standard (recommandé) \\
\hline
Q & 25\% & Documents photocopiés, archivage long terme \\
\hline
H & 30\% & Environnements difficiles, manipulations fréquentes \\
\hline
\end{tabular}
\caption{Niveaux de correction d'erreur des QR codes}
\end{table}

\subsection{Glossaire}

\begin{description}[leftmargin=5cm, style=nextline]
    \item[AES] Advanced Encryption Standard - Algorithme de chiffrement symétrique

    \item[Attestation de réussite] Document officiel certifiant qu'un étudiant a validé un semestre (moyenne $\geq$ 10/20)

    \item[Base de notation] Échelle sur laquelle les notes sont données (10, 20, ou 100)

    \item[Chiffrement compact] Méthode de chiffrement produisant une sortie de taille réduite, adaptée aux QR codes

    \item[Coefficient] Poids d'un EC dans le calcul de la moyenne d'une UE

    \item[Configuration de classe] Ensemble structuré définissant les semestres, UEs et ECs d'une classe

    \item[Crédits ECTS] European Credit Transfer and Accumulation System - Système européen de crédits académiques

    \item[EC] Élément Constitutif - Matière ou cours au sein d'une UE

    \item[Excel] Format de fichier tableur (.xlsx, .xls) utilisé pour importer les données

    \item[Mapping] Association entre les colonnes Excel et les ECs de la configuration

    \item[Matricule] Numéro d'identification unique d'un étudiant

    \item[Mention] Appréciation qualitative basée sur la moyenne (Très Bien, Bien, Assez Bien, Passable, Ajourné)

    \item[Mode démo] Mode de fonctionnement sans licence, avec restrictions

    \item[Multi-semestres] Fonctionnalité permettant de générer des documents pour plusieurs semestres simultanément

    \item[OOM] Out Of Memory - Erreur survenant quand la mémoire RAM est insuffisante

    \item[PDF fusionné] Un seul fichier PDF contenant plusieurs documents assemblés séquentiellement

    \item[Préréglage] Thème prédéfini prêt à l'emploi

    \item[QR code] Quick Response Code - Code-barres 2D lisible par smartphone

    \item[Relevé de notes] Document officiel récapitulant les notes, moyennes et crédits d'un étudiant pour un semestre

    \item[Semestre] Période académique (généralement 6 mois) regroupant plusieurs UEs

    \item[Thème] Ensemble de paramètres visuels (couleurs, polices, mise en page) appliqué aux documents

    \item[UE] Unité d'Enseignement - Regroupement cohérent d'ECs valant un certain nombre de crédits

    \item[ZIP] Format d'archive compressée permettant de regrouper plusieurs fichiers
\end{description}

\subsection{Limites techniques}

\subsubsection{Limites pour les relevés de notes}

\begin{itemize}[leftmargin=*]
    \item Nombre maximum de relevés par génération : \textbf{500} (recommandé : 200)
    \item Nombre maximum de semestres par configuration : \textbf{12}
    \item Nombre maximum d'UEs par semestre : \textbf{20}
    \item Nombre maximum d'ECs par UE : \textbf{15}
    \item Taille maximale du fichier Excel : \textbf{50 Mo}
    \item Nombre maximum de lignes Excel : \textbf{10 000}
\end{itemize}

\subsubsection{Limites pour les attestations}

\begin{itemize}[leftmargin=*]
    \item Nombre maximum d'attestations par génération : \textbf{1000} (recommandé : 500)
    \item Nombre maximum de thèmes personnalisés sauvegardés : \textbf{50}
    \item Taille maximale du logo : \textbf{2 Mo} (format PNG ou JPG)
    \item Résolution logo recommandée : \textbf{800x200 pixels}
\end{itemize}

\subsubsection{Limites pour les QR codes}

\begin{itemize}[leftmargin=*]
    \item Nombre maximum de PDFs traités simultanément : \textbf{200}
    \item Taille maximale d'un PDF individuel : \textbf{50 Mo}
    \item Taille QR code : Min \textbf{40 px}, Max \textbf{200 px}
    \item Nombre de colonnes dans QR : Recommandé max \textbf{6}
    \item Longueur totale contenu QR : Max \textbf{2000 caractères}
\end{itemize}

\subsection{Compatibilité des formats}

\subsubsection{Formats de fichiers en entrée}

\begin{table}[h]
\centering
\begin{tabular}{|l|p{5cm}|p{5cm}|}
\hline
\textbf{Type} & \textbf{Formats acceptés} & \textbf{Notes} \\
\hline
Données étudiants & .xlsx, .xls, .csv & Excel 2007+ recommandé \\
\hline
Documents PDF & .pdf & Version PDF 1.4 à 1.7 \\
\hline
Images (logo) & .png, .jpg, .jpeg & PNG avec transparence recommandé \\
\hline
Configurations & .json & Export de l'application \\
\hline
Thèmes & .json & Export de l'application \\
\hline
\end{tabular}
\caption{Formats de fichiers supportés en entrée}
\end{table}

\subsubsection{Formats de fichiers en sortie}

\begin{table}[h]
\centering
\begin{tabular}{|l|p{10cm}|}
\hline
\textbf{Format} & \textbf{Description} \\
\hline
PDF & Documents générés (relevés, attestations) au format PDF/A \\
\hline
ZIP & Archives compressées (compression DEFLATE) \\
\hline
JSON & Exports de configurations, thèmes, historique \\
\hline
\end{tabular}
\caption{Formats de fichiers générés}
\end{table}

\subsection{FAQ - Foire aux questions}

\subsubsection{Questions générales}

\textbf{Q : L'application nécessite-t-elle une connexion Internet ?}

R : Non. L'application fonctionne entièrement hors ligne. Seule la vérification des mises à jour nécessite Internet (optionnel).

\textbf{Q : Mes données sont-elles sécurisées ?}

R : Oui. Toutes les données restent sur votre ordinateur local. Rien n'est envoyé sur Internet.

\textbf{Q : Puis-je utiliser l'application sur plusieurs ordinateurs ?}

R : Oui, mais vous devez installer l'application sur chaque ordinateur. Les configurations peuvent être partagées via export/import.

\textbf{Q : L'application est-elle compatible avec macOS et Linux ?}

R : Oui. Des versions sont disponibles pour Windows, macOS et Linux.

\subsubsection{Questions sur les relevés}

\textbf{Q : Comment ajouter une nouvelle UE à une configuration existante ?}

R : Menu \menu{Config des relevés} > Sélectionnez votre configuration > Cliquez sur le semestre > \bouton{Ajouter une UE}.

\textbf{Q : Les moyennes sont arrondies comment ?}

R : Deux décimales. Exemple : 12.456 devient 12.46.

\textbf{Q : Que se passe-t-il si une note est manquante ?}

R : La note est ignorée dans le calcul. Seules les notes présentes sont prises en compte.

\textbf{Q : Puis-je changer le barème des mentions ?}

R : Non, le barème est fixe : Très Bien ($\geq$16), Bien (14-15.99), Assez Bien (12-13.99), Passable (10-11.99), Ajourné (<10).

\subsubsection{Questions sur les attestations}

\textbf{Q : Pourquoi un étudiant avec moyenne 9.9 ne reçoit pas d'attestation ?}

R : Les attestations de réussite sont réservées aux étudiants admis (moyenne $\geq$ 10/20). C'est un filtrage automatique.

\textbf{Q : Puis-je forcer la génération d'une attestation pour un étudiant ajourné ?}

R : Non. L'application n'autorise que les étudiants ayant validé leur semestre.

\textbf{Q : Combien de thèmes puis-je créer ?}

R : Maximum 50 thèmes personnalisés sauvegardés.

\textbf{Q : Les attestations bilingues nécessitent-elles deux générations ?}

R : Non. Une seule génération. Sélectionnez la langue dans les paramètres, l'application utilise les bonnes colonnes Excel.

\subsubsection{Questions sur les QR codes}

\textbf{Q : Le QR code est-il ajouté sur toutes les pages ou seulement la première ?}

R : Sur \textbf{TOUTES les pages} de chaque document. C'est une caractéristique unique de cette application.

\textbf{Q : Puis-je scanner un QR code directement depuis l'écran ?}

R : Oui. Ouvrez le PDF sur votre ordinateur et scannez avec un smartphone.

\textbf{Q : Quelle application de scan recommandez-vous ?}

R : iOS : Appareil photo natif. Android : Google Lens. Ou toute app "QR Code Reader".

\textbf{Q : Le QR code chiffré peut-il être déchiffré ?}

R : Uniquement avec la clé appropriée (basée sur le matricule de l'étudiant).

\subsection{Support et ressources}

\subsubsection{Canaux de support}

\begin{itemize}[leftmargin=*]
    \item \textbf{Email} : support@fmsp.edu
    \item \textbf{Documentation en ligne} : https://docs.fmsp.edu/releves
    \item \textbf{Système de tickets} : https://support.fmsp.edu
    \item \textbf{Menu Aide} : Intégré dans l'application (F1)
\end{itemize}

\subsubsection{Ressources utiles}

\begin{itemize}[leftmargin=*]
    \item Guide de démarrage rapide (PDF) : Disponible dans le menu \menu{Aide}
    \item Tutoriels vidéo : https://videos.fmsp.edu/tutoriels
    \item Templates Excel : https://download.fmsp.edu/templates
    \item Forum communautaire : https://forum.fmsp.edu
\end{itemize}

\subsubsection{Signaler un bug}

Pour signaler un problème technique :

\begin{enumerate}[leftmargin=*]
    \item Menu \menu{Aide} > \bouton{Rapporter un problème}
    \item Ou envoyez un email à bugs@fmsp.edu
    \item Incluez :
    \begin{itemize}
        \item Version de l'application
        \item Description détaillée
        \item Étapes pour reproduire
        \item Captures d'écran
        \item Fichiers de logs (si applicable)
    \end{itemize}
\end{enumerate}

\subsection{Historique des versions}

\subsubsection{Version 2.0 (Actuelle)}

\textbf{Nouvelles fonctionnalités :}
\begin{itemize}[leftmargin=*]
    \item Génération multi-semestres améliorée
    \item Support bilingue FR/EN pour attestations
    \item 6 nouveaux préréglages de thèmes
    \item QR codes sur toutes les pages des PDFs multi-pages
    \item Export de configurations et thèmes
    \item Historique des documents générés
\end{itemize}

\textbf{Améliorations :}
\begin{itemize}[leftmargin=*]
    \item Optimisation mémoire pour générations massives
    \item Interface utilisateur repensée
    \item Performances de génération PDF améliorées de 30\%
    \item Meilleure gestion des erreurs Excel
\end{itemize}

\subsubsection{Version 1.5}

\textbf{Fonctionnalités :}
\begin{itemize}[leftmargin=*]
    \item Première version avec QR codes chiffrés
    \item Thèmes personnalisables
    \item Export ZIP avec compression
\end{itemize}

\subsubsection{Version 1.0 (Initiale)}

\textbf{Fonctionnalités de base :}
\begin{itemize}[leftmargin=*]
    \item Génération de relevés de notes
    \item Génération d'attestations de réussite
    \item Configuration académique basique
    \item Export PDF simple
\end{itemize}

\vfill

\begin{center}
\rule{\textwidth}{0.4pt}

\vspace{0.5cm}

\textit{Fin du guide d'utilisation complet}

\vspace{0.3cm}

\textbf{Version 2.0 - Corrigée et détaillée}

\vspace{0.2cm}

Basée sur l'analyse complète du code source

\vspace{0.3cm}

\today

\vspace{0.5cm}

\textit{Faculté de Médecine et Sciences Pharmaceutiques}

\vspace{0.2cm}

Pour toute question : support@fmsp.edu

\vspace{0.5cm}

\rule{\textwidth}{0.4pt}
\end{center}

\end{document}


\end{document}
